\documentclass[a4paper,12pt]{article}

\usepackage[T1]{fontenc}
\usepackage[italian]{babel}
\usepackage{graphicx}
\usepackage{geometry}
\usepackage{tabularx}
\usepackage[table]{xcolor}
\usepackage[most]{tcolorbox}
\usepackage{float}
\usepackage{fancyhdr}
\usepackage[hidelinks]{hyperref}
\usepackage{amsmath}
\usepackage{longtable}
\usepackage{eurosym}
\usepackage{pgfplots}
\pgfplotsset{compat=1.18}
\usetikzlibrary{patterns}
\graphicspath{{./}{../../}}

\newcommand{\groupmail}{alittlebyte.swe@gmail.com}
\newcommand{\grouplogo}{../../.github/site-src/assets/logo_alittlebyte.png}
\newcommand{\groupsymbol}{../../.github/site-src/assets/solo_logo.png}


\geometry{
    left=2.5cm,
    right=2.5cm,
    top=2.5cm,
    bottom=2.5cm,
    headheight=331pt,
    includeheadfoot
}

\pagestyle{fancy}
\fancyhf{}

\renewcommand{\headrulewidth}{0pt}
\fancyhead{}

\renewcommand{\footrulewidth}{0.4pt}
\fancyfoot[L]{\includegraphics[height=0.6cm]{\groupsymbol}\hspace{0.45em}aLittleByte}
\fancyfoot[R]{Piano di Qualifica}

\fancypagestyle{firstpage}{
    \fancyhf{}
    \renewcommand{\headrulewidth}{0pt}
    \renewcommand{\footrulewidth}{0.4pt}
    \fancyhead[C]{%
        \makebox[\textwidth][c]{%
            \begin{tabular}{c}
                \includegraphics[height=10cm]{\grouplogo}\\[1ex]
                \small\texttt{\groupmail}
            \end{tabular}%
        }
    }
    \fancyfoot[L]{\includegraphics[height=0.6cm]{\groupsymbol}\hspace{0.45em}aLittleByte}
    \fancyfoot[R]{Piano di Qualifica}

    \setcounter{secnumdepth}{4}
    \setcounter{tocdepth}{4}
}

\providecommand{\glslink}[2]{\href{http://127.0.0.1:8765/glossary-html\#gls-#1}{\underline{#2}\textsuperscript{\scriptsize G}}}

\begin{document}

\thispagestyle{firstpage}

\vspace*{0.5cm}

\begin{center}
    {\LARGE \textbf{Piano di Qualifica}}
\end{center}

\vspace{1.5cm}

\begin{center}
    \renewcommand{\arraystretch}{1.3}
    \begin{tcolorbox}[
        width=0.92\textwidth,
        colback=gray!6,
        colframe=gray!45,
        boxrule=0.5pt,
        arc=2mm,
        boxsep=0pt,
        left=8pt, right=8pt, top=8pt, bottom=8pt
    ]
        \begin{tabularx}{\linewidth}{>{\bfseries}p{3.5cm} X}
            Gruppo      & aLittleByte (gruppo 19) \\
            Redattori   & Eleonora Bellet, Diego Belluz \\
            Verificatori& Leonardo Soligo, Angelica Gastaldello\\
            Destinatari & Prof. Tullio Vardanega, Prof. Riccardo Cardin \\
            Versione    & 0.6.0 \\
            Data        & 28/05/2026 \\
        \end{tabularx}
    \end{tcolorbox}
\end{center}

\newpage

\newgeometry{left=2.5cm,right=2.5cm,top=2.5cm,bottom=2.5cm,includefoot}

\begin{center}
    {\Large \textbf{Registro delle Modifiche}}
\end{center}

\vspace{0.35cm}

\renewcommand{\arraystretch}{1.35}
\rowcolors{2}{gray!8}{white}
\begin{center}
    {\footnotesize
    \begin{tabularx}{\textwidth}{|p{1.4cm}|p{1.8cm}|X|p{2.4cm}|p{2.3cm}|p{2.3cm}|}
        \hline
        \rowcolor{gray!20}
        \textbf{Versione} & \textbf{Data} & \textbf{Descrizione} & \textbf{Redattore} & \textbf{Verificatore} & \textbf{Approvatore} \\
        \hline
        0.6.0 & 28/05/2026 & Aggiunta sezione 6 - Iniziative di miglioramento. & - & - & - \\
        \hline
        0.5.2 & 28/05/2026 & Correzione applicabilità metriche MPC-CC e MPC-TSR. & - & - & - \\
        \hline
        0.5.1 & 17/05/2026 & Aggiornamento metriche RSI, PRI, TE, LE e ricalcolo QMS. & - & - & - \\
        \hline
        0.5.0 & 17/05/2026 & Aggiunta sezione 5 - Cruscotto di qualità & - & - & - \\
        \hline
        0.4.3 & 15/05/2026 & Test di accettazione e tabelle di tracciamento & - & - & - \\
        \hline
        0.4.2 & 14/05/2026 & Test di sistema - qualità e vincolo & - & - & - \\
        \hline
        0.4.1 & 13/05/2026 & Test di sistema - funzionali & - & - & - \\
        \hline
        0.4.0 & 12/05/2026 & Aggiunta sezione 4 - Metodi di testing & - & - & - \\
        \hline
        0.3.2 & 09/05/2026 & Correzioni soglie metriche MPC e MPD & - & - & - \\
        \hline
        0.3.1 & 08/05/2026 & Aggiunta metriche MPD per usabilità, mantenibilità e portabilità & - & - & - \\
        \hline
        0.3.0 & 07/05/2026 & Aggiunta sezione 3 - Qualità di prodotto & - & - & - \\
        \hline
        0.2.1 & 05/05/2026 & Aggiunta metriche MPC per processi organizzativi & - & - & - \\
        \hline
        0.2.0 & 03/05/2026 & Aggiunta sezione 2 - Qualità di processo & - & - & - \\
        \hline
        0.1.1 & 01/05/2026 & Completamento riferimenti normativi e informativi & - & - & - \\
        \hline
        0.1.0 & 29/04/2026 & Prima stesura del documento & - & - & - \\
        \hline
    \end{tabularx}
    }

\end{center}
\newpage

\begin{center} 
    {\Large \textbf{Indice}}
\end{center}

\vspace{0.6cm}

\makeatletter
\@starttoc{toc}
\makeatother

\newpage
\pagenumbering{arabic}
\setcounter{page}{4}
\fancyfoot[R]{Piano di Qualifica\hspace{0.8em}\thepage}

\section{Introduzione}

\subsection{Finalità del documento}
Il presente documento, denominato \textit{\glslink{piano-di-qualifica}{Piano di Qualifica}}, ha lo scopo di definire le strategie, le procedure e le metriche adottate dal gruppo \textit{aLittleByte} per garantire la qualità del prodotto software e dei processi produttivi relativi al progetto C5 (\glslink{nexum}{NEXUM}), proposto dall'azienda \textit{Eggon}.

In particolare, questo documento si prefigge di:
\begin{itemize}
    \item \textbf{Definire gli obiettivi di qualità:} specificare i target qualitativi per il processo di sviluppo (efficienza, stabilità) e per il prodotto software (funzionalità, affidabilità, manutenibilità), in conformità con gli standard ISO/IEC 12207 e ISO/IEC 9126;
    \item \textbf{Identificare le metriche:} selezionare gli indicatori quantitativi più idonei per monitorare il raggiungimento degli obiettivi, fissando per ciascuno le soglie di accettazione e di ottimalità;
    \item \textbf{Pianificare le attività di verifica e validazione:} descrivere le metodologie di test (unità, integrazione, sistema, accettazione) e le procedure di \glslink{analisi-statica}{analisi statica} del codice e della documentazione;
    \item \textbf{Monitorare l'andamento del progetto:} fornire un resoconto puntuale delle misurazioni effettuate durante le varie fasi del ciclo di vita, permettendo al team di individuare tempestivamente criticità e attuare azioni correttive.
\end{itemize}

\subsection{Glossario}
Al fine di evitare ambiguità e garantire una comprensione uniforme della terminologia utilizzata, è stato redatto un documento esterno denominato \textit{\glslink{glossario}{Glossario}}.
I termini tecnici, gli acronimi e le parole con un significato specifico all'interno del progetto sono contrassegnati nel testo da una ``G'' in pedice (es. parola\textsubscript{\scalebox{0.6}{\textbf{G}}}). La loro definizione completa è consultabile nel \textit{\glslink{glossario}{Glossario}}.

\subsection{Riferimenti}
\subsubsection{Riferimenti normativi}
\begin{itemize}
    \item \textbf{\glslink{capitolato}{Capitolato} d'appalto C5 - \glslink{nexum}{NEXUM} (Eggon):}\\
    \url{https://www.math.unipd.it/~tullio/IS-1/2025/Progetto/C5.pdf}

    \item \textbf{\glslink{norme-di-progetto}{Norme di Progetto}:}\\
    Documento interno del gruppo \textit{aLittleByte} che definisce le regole, i ruoli e le procedure operative.\\
    \url{https://github.com/aLittleByte-19/Documentazione}

    \item \textbf{Regolamento del progetto didattico:}\\
    \url{https://www.math.unipd.it/~tullio/IS-1/2025/Dispense/PD1.pdf}

    \item \textbf{\glslink{piano-di-progetto}{Piano di Progetto}:}\\
    Documento interno del gruppo \textit{aLittleByte}.\\
    \url{https://github.com/aLittleByte-19/Documentazione}
\end{itemize}

\subsubsection{Riferimenti informativi}
\begin{itemize}
    \item \textbf{\glslink{glossario}{Glossario}:}\\
    Documento interno del gruppo \textit{aLittleByte} contenente le definizioni dei termini tecnici.\\
    \url{https://github.com/aLittleByte-19/Documentazione}

    \item \textbf{Standard ISO/IEC 12207:1995:}\\
    \textit{Information technology -- Software life cycle processes}.\\
    \url{https://www.math.unipd.it/~tullio/IS-1/2009/Approfondimenti/ISO_12207-1995.pdf}

    \item \textbf{Standard ISO/IEC 9126:}\\
    \textit{Software engineering -- Product quality}.\\
    \url{https://it.wikipedia.org/wiki/ISO/IEC_9126}

    \item \textbf{Slide T07 - Qualità del software:}\\
    \url{https://www.math.unipd.it/~tullio/IS-1/2025/Dispense/T07.pdf}\\

    \item \textbf{Slide T08 - Qualità del software:}\\
    \url{https://www.math.unipd.it/~tullio/IS-1/2025/Dispense/T08.pdf}\\
\end{itemize}

\section{Qualità di processo}
La qualità di processo descrive quanto le attività di sviluppo siano pianificate, controllate e migliorate nel tempo. Il riferimento metodologico adottato è lo standard ISO/IEC 12207, adattato al \glslink{way-of-working-wow}{Way of Working} del gruppo.
Le metriche di qualità del processo misurano l'efficacia, l'efficienza e il controllo delle attività necessarie per sviluppare, gestire e consegnare il prodotto software. Il loro scopo è di monitorare l'aderenza alla \glslink{pianificazione}{pianificazione}, la stabilità dei processi, la sostenibilità dei costi e il miglioramento costante dei risultati. Queste metriche non valutano il prodotto finito, ma il modo in cui il prodotto viene realizzato.
Il monitoraggio continuo consente di confrontare i valori rilevati con le soglie stabilite, garantendo trasparenza verso i committenti e coerenza interna nelle scelte operative.

In questo documento, tali misure vengono identificate tramite la sigla \textbf{MPC} (Metriche di Processo e Controllo). Questo identificativo permette di classificare e tracciare tutte le misurazioni relative alla gestione dei costi, all'avanzamento temporale, alla qualità della documentazione e all'efficienza dei processi interni al team di sviluppo.

\subsection{Processi primari}
\subsubsection{Fornitura}
\begin{table}[H]
    \centering
    \renewcommand{\arraystretch}{1.2}
    \rowcolors{2}{gray!8}{white}
    \begin{tabularx}{\textwidth}{|p{2.4cm}|X|p{3cm}|p{3cm}|}
        \hline
        \rowcolor{gray!20}
        \textbf{ID} & \textbf{Nome Metrica} & \textbf{Accettabile} & \textbf{Ottimo} \\
        \hline
        MPC-EV & Earned Value & $\geq 0$ & $\leq$ EAC \\
        MPC-PV & Planned Value & $\geq 0$ & $\leq$ BAC \\
        MPC-AC & Actual Cost & $\geq 0$ & $\leq$ EAC \\
        MPC-CPI & Cost Performance Index & $\geq 0,9$ & $1$ \\
        MPC-SPI & Schedule Performance Index & $\geq 0,9$ & $1$ \\
        MPC-EAC & Estimated At Completion & $\pm 5\%$ rispetto al BAC & BAC \\
        MPC-ETC & Estimate To Complete & $\geq 0$ & $\leq$ EAC \\
        MPC-TEAC & Time Estimate At Completion & $\geq 0$ & $\leq$ Durata pianificata \\
        \hline
    \end{tabularx}
    \caption{Valori per misurare la qualità della fornitura.}
\end{table}

\subsubsection{Sviluppo}
\begin{table}[H]
    \centering
    \renewcommand{\arraystretch}{1.2}
    \rowcolors{2}{gray!8}{white}
    \begin{tabularx}{\textwidth}{|p{2.4cm}|X|p{3cm}|p{3cm}|}
        \hline
        \rowcolor{gray!20}
        \textbf{ID} & \textbf{Nome Metrica} & \textbf{Accettabile} & \textbf{Ottimo} \\
        \hline
        MPC-RSI & Requirements Stability Index & $\geq 70\%$ & $100\%$ \\
        \hline
    \end{tabularx}
    \caption{Valori per misurare la qualità dello sviluppo.}
\end{table}

\subsection{Processi di supporto}
\subsubsection{Documentazione}
\begin{table}[H]
    \centering
    \renewcommand{\arraystretch}{1.2}
    \rowcolors{2}{gray!8}{white}
    \begin{tabularx}{\textwidth}{|p{2.4cm}|X|p{3cm}|p{3cm}|}
        \hline
        \rowcolor{gray!20}
        \textbf{ID} & \textbf{Nome Metrica} & \textbf{Accettabile} & \textbf{Ottimo} \\
        \hline
        \glslink{indice-di-gulpease-mpc-ig}{MPC-IG} & \glslink{indice-di-gulpease-mpc-ig}{Indice di Gulpease} del documento & $\geq 60$ & $\geq 80$ \\
        MPC-CO & Correttezza ortografica & $0$ errori & $0$ errori \\
        \hline
    \end{tabularx}
    \caption{Valori per misurare la qualità della documentazione.}
\end{table}

\subsubsection{Verifica}
\begin{table}[H]
    \centering
    \renewcommand{\arraystretch}{1.2}
    \rowcolors{2}{gray!8}{white}
    \begin{tabularx}{\textwidth}{|p{2.4cm}|X|p{3cm}|p{3cm}|}
        \hline
        \rowcolor{gray!20}
        \textbf{ID} & \textbf{Nome Metrica} & \textbf{Accettabile} & \textbf{Ottimo} \\
        \hline
        MPC-CC & Code Coverage & $\geq 80\%$ & $100\%$ \\
        MPC-TSR & Test Success Rate & $100\%$ & $100\%$ \\
        \hline
    \end{tabularx}
    \caption{Valori per misurare la qualità della verifica.}
\end{table}

\subsubsection{Gestione della qualità}
\begin{table}[H]
    \centering
    \renewcommand{\arraystretch}{1.2}
    \rowcolors{2}{gray!8}{white}
    \begin{tabularx}{\textwidth}{|p{2.4cm}|X|p{3cm}|p{3cm}|}
        \hline
        \rowcolor{gray!20}
        \textbf{ID} & \textbf{Nome Metrica} & \textbf{Accettabile} & \textbf{Ottimo} \\
        \hline
        MPC-QMS & Quality Metrics Satisfied & $\geq 80\%$ & $100\%$ \\
        \hline
    \end{tabularx}
    \caption{Valori per misurare la qualità della gestione della qualità.}
\end{table}

\subsection{Processi organizzativi}
\subsubsection{Gestione dei processi}
\begin{table}[H]
    \centering
    \renewcommand{\arraystretch}{1.2}
    \rowcolors{2}{gray!8}{white}
    \begin{tabularx}{\textwidth}{|p{2.4cm}|X|p{3cm}|p{3cm}|}
        \hline
        \rowcolor{gray!20}
        \textbf{ID} & \textbf{Nome Metrica} & \textbf{Accettabile} & \textbf{Ottimo} \\
        \hline
        MPC-TE & Time Efficiency & $\geq 50\%$ & $100\%$ \\
        MPC-PRI & Percentuale Rischi Inattesi & $\leq 20\%$ & $0\%$ \\
        MPC-LE & Labor Efficiency & $\geq 70\%$ & $\geq 100\%$ \\
        \hline
    \end{tabularx}
    \caption{Valori per misurare la qualità della gestione dei processi.}
\end{table}

\section{Qualità di prodotto}
La qualità di prodotto misura le proprietà interne ed esterne del software finale. Lo standard di riferimento è ISO/IEC 9126, adottato come base per definire le metriche di prodotto.
Tali metriche consentono di verificare che il sistema rispetti i requisiti e che l'esperienza d'uso risulti stabile e prevedibile nelle condizioni operative attese.

In questo documento, tali misure vengono identificate tramite la sigla \textbf{MPD} (Metriche di Prodotto). Questo identificativo permette di classificare e monitorare le caratteristiche del software, facilitando la verifica del raggiungimento degli obiettivi qualitativi prefissati.

\subsection{Funzionalità}
\begin{table}[H]
    \centering
    \renewcommand{\arraystretch}{1.2}
    \rowcolors{2}{gray!8}{white}
    \begin{tabularx}{\textwidth}{|p{2.4cm}|X|p{3cm}|p{3cm}|}
        \hline
        \rowcolor{gray!20}
        \textbf{ID} & \textbf{Nome Metrica} & \textbf{Accettabile} & \textbf{Ottimo} \\
        \hline
        MPD01 & Requisiti obbligatori soddisfatti & $100\%$ & $100\%$ \\
        MPD02 & Requisiti opzionali soddisfatti & $0\%$ & $100\%$ \\
        MPD03 & \glslink{intelligenza-artificiale-ai}{AI} Acceptance Rate (\glslink{rating}{rating} $\geq 3/5$) & $\geq 60\%$ & $\geq 80\%$ \\
        \hline
    \end{tabularx}
    \caption{Valori per misurare la qualità delle funzionalità.}
\end{table}

\subsection{Affidabilità}
\begin{table}[H]
    \centering
    \renewcommand{\arraystretch}{1.2}
    \rowcolors{2}{gray!8}{white}
    \begin{tabularx}{\textwidth}{|p{2.4cm}|X|p{3cm}|p{3cm}|}
        \hline
        \rowcolor{gray!20}
        \textbf{ID} & \textbf{Nome Metrica} & \textbf{Accettabile} & \textbf{Ottimo} \\
        \hline
        MPD04 & \glslink{branch}{Branch} Coverage & $\geq 70\%$ & $\geq 85\%$ \\
        MPD05 & Defect Density & $\leq 3$ / \glslink{kloc}{KLOC} & $\leq 1$ / \glslink{kloc}{KLOC} \\
        \hline
    \end{tabularx}
    \caption{Valori per misurare la qualità dell'affidabilità.}
\end{table}

\subsection{Efficienza}
\begin{table}[H]
    \centering
    \renewcommand{\arraystretch}{1.2}
    \rowcolors{2}{gray!8}{white}
    \begin{tabularx}{\textwidth}{|p{2.4cm}|X|p{3cm}|p{3cm}|}
        \hline
        \rowcolor{gray!20}
        \textbf{ID} & \textbf{Nome Metrica} & \textbf{Accettabile} & \textbf{Ottimo} \\
        \hline
        MPD06 & \glslink{latenza}{Latenza} interfaccia utente & $\leq 2$ sec & $\leq 0,5$ sec \\
        MPD07 & \glslink{latenza}{Latenza} generazione testo \glslink{intelligenza-artificiale-ai}{AI} Assistant & $\leq 5$ sec & $\leq 3$ sec \\
        MPD08 & \glslink{latenza}{Latenza} generazione immagine \glslink{intelligenza-artificiale-ai}{AI} Assistant & $\leq 10$ sec & $\leq 5$ sec \\
        MPD09 & \glslink{latenza}{Latenza} analisi documentale \glslink{intelligenza-artificiale-ai}{AI} Co-Pilot & $\leq 10$ sec & $\leq 5$ sec \\
        \hline
    \end{tabularx}
    \caption{Valori per misurare la qualità dell'efficienza.}
\end{table}

\subsection{Usabilità}
\begin{table}[H]
    \centering
    \renewcommand{\arraystretch}{1.2}
    \rowcolors{2}{gray!8}{white}
    \begin{tabularx}{\textwidth}{|p{2.4cm}|X|p{3cm}|p{3cm}|}
        \hline
        \rowcolor{gray!20}
        \textbf{ID} & \textbf{Nome Metrica} & \textbf{Accettabile} & \textbf{Ottimo} \\
        \hline
        MPD10 & Profondità Massima di Navigazione & $\leq 5$ livelli & $\leq 3$ livelli \\
        MPD11 & Learning Time (funzioni principali) & $\leq 30$ min & $\leq 15$ min \\
        \hline
    \end{tabularx}
    \caption{Valori per misurare la qualità dell'usabilità.}
\end{table}

\subsection{Mantenibilità}
\begin{table}[H]
    \centering
    \renewcommand{\arraystretch}{1.2}
    \rowcolors{2}{gray!8}{white}
    \begin{tabularx}{\textwidth}{|p{2.4cm}|X|p{3cm}|p{3cm}|}
        \hline
        \rowcolor{gray!20}
        \textbf{ID} & \textbf{Nome Metrica} & \textbf{Accettabile} & \textbf{Ottimo} \\
        \hline
        MPD12 & Code Smell & $\leq 5$ & $0$ \\
        MPD13 & Coefficient of Coupling & $\leq 0{,}7$ & $\leq 0{,}4$ \\
        MPD14 & Cyclomatic Complexity (per metodo) & $\leq 15$ & $\leq 10$ \\
        \hline
    \end{tabularx}
    \caption{Valori per misurare la qualità della mantenibilità.}
\end{table}

\subsection{Portabilità}
\begin{table}[H]
    \centering
    \renewcommand{\arraystretch}{1.2}
    \rowcolors{2}{gray!8}{white}
    \begin{tabularx}{\textwidth}{|p{2.4cm}|X|p{3cm}|p{3cm}|}
        \hline
        \rowcolor{gray!20}
        \textbf{ID} & \textbf{Nome Metrica} & \textbf{Accettabile} & \textbf{Ottimo} \\
        \hline
        MPD15 & Compatibilità Cross-Browser & $\geq 80\%$ & $100\%$ \\
        \hline
    \end{tabularx}
    \caption{Valori per misurare la qualità della portabilità.}
\end{table}

\section{Metodi di testing}
La presente sezione descrive le attività di testing adottate nel progetto e le metriche utilizzate per valutare l'efficacia del processo di verifica.
Le attività di testing forniscono evidenza oggettiva del corretto funzionamento dell'intero sistema e supportano la valutazione delle metriche di qualità del prodotto discusse in questo documento.

I test avranno un codice univoco, la descrizione, il requisito di riferimento e lo stato attuale. I codici per descrivere gli stati dei test sono i seguenti:
\begin{itemize}
    \item \textbf{NI}: Non Implementato
    \item \textbf{I}: Implementato, ma non ancora verificato
    \item \textbf{S}: Superato, ovvero il test è stato eseguito ed ha restituito un esito positivo
    \item \textbf{NS}: Non Superato, ovvero il test è stato eseguito ma ha restituito un esito negativo
\end{itemize}

\subsection{Copertura del codice}
La copertura del codice (Code Coverage) misura la percentuale di codice sorgente eseguita durante l'esecuzione dei test automatici. Tale metrica consente di valutare il grado di verifica del software ed è direttamente collegata alla metrica \textbf{MPC-CC}.
Il valore minimo accettabile è fissato all'80\%.

\subsection{Test unitari}
I test unitari verificano il corretto funzionamento delle singole unità software in isolamento. L'attenzione è rivolta alle componenti critiche del backend: logica di elaborazione dei \glslink{prompt}{prompt}, validazione degli input, gestione dei \glslink{metadati}{metadati} documentali e algoritmi di classificazione \glslink{intelligenza-artificiale-ai}{AI}. I test sono realizzati con il framework \glslink{pest}{Pest} e contribuiscono al miglioramento delle metriche \textbf{MPC-TSR} e \textbf{MPC-CC}.

\subsection{Test di integrazione}
I \glslink{test-di-integrazione}{test di integrazione} verificano il corretto comportamento delle interazioni tra i componenti del sistema. Considerata la natura distribuita dell'architettura, tali test risultano fondamentali per il raggiungimento di un solido risultato.

\renewcommand{\arraystretch}{1.15}
\rowcolors{2}{gray!8}{white}
\begin{longtable}{|p{2.2cm}|p{8.5cm}|p{3.8cm}|p{1.5cm}|}
    \caption{Test di integrazione} \label{tab:ti} \\
    \hline
    \rowcolor{gray!20}
    \textbf{Codice} & \textbf{Descrizione} & \textbf{Moduli} & \textbf{Stato} \\
    \hline
    \endfirsthead
    \hline
    \rowcolor{gray!20}
    \textbf{Codice} & \textbf{Descrizione} & \textbf{Moduli} & \textbf{Stato} \\
    \hline
    \endhead
    \hline
    \endfoot
    \hline
    \endlastfoot
    TI-001 & Verifica che il sistema scambi correttamente i dati tra frontend e backend tramite API REST. & UI (Blade) $\leftrightarrow$ \glslink{laravel}{Laravel} 12 & NI \\
    \hline
    TI-002 & Verifica che il sistema invii il \glslink{prompt}{prompt} e riceva la risposta generata dal servizio \glslink{intelligenza-artificiale-ai}{AI}. & Backend $\leftrightarrow$ \glslink{amazon-bedrock}{Amazon Bedrock} & NI \\
    \hline
    TI-003 & Verifica che il sistema persista e recuperi correttamente i dati applicativi dal database. & Backend $\leftrightarrow$ \glslink{postgresql}{PostgreSQL} & NI \\
    \hline
    TI-004 & Verifica che il sistema esegua la pipeline \glslink{ocr-riconoscimento-ottico-dei-caratteri}{OCR} e la classificazione dei documenti tramite \glslink{ai-analyst}{AI Analyst}. & Backend $\leftrightarrow$ \glslink{ai-analyst}{AI Analyst} & NI \\
    \hline
    TI-005 & Verifica che il sistema gestisca correttamente l'upload, lo storage e l'accesso \glslink{intelligenza-artificiale-ai}{ai} file. & Backend $\leftrightarrow$ \glslink{minio}{MinIO} (S3) & NI \\
    \hline
    TI-006 & Verifica che il sistema esegua correttamente i job in modalità asincrona tramite le code. & Backend $\leftrightarrow$ \glslink{redis}{Redis} Queue & NI \\
    \hline
\end{longtable}

\subsection{Test di regressione}
I \glslink{test-di-regressione}{test di regressione} vengono eseguiti al termine di ogni \glslink{sprint}{sprint}, dopo l'integrazione di nuove funzionalità o la correzione di difetti, con lo scopo di accertare che il comportamento già verificato non sia stato compromesso. Consistono nella riesecuzione automatica dell'intera suite di test unitari e di integrazione preesistente. Un aumento dei fallimenti costituisce un indicatore di instabilità che richiede analisi e risoluzione prima del rilascio.

\subsection{Test di sistema}
I \glslink{test-di-sistema}{test di sistema} verificano il corretto comportamento complessivo dell'applicazione in un ambiente il più possibile simile a quello di utilizzo reale. La tabella seguente raccoglie in modo unitario tutti i \glslink{test-di-sistema}{test di sistema}, raggruppati per tipologia di requisito: funzionali, di qualità e di vincolo.

\renewcommand{\arraystretch}{1.15}
\rowcolors{2}{gray!8}{white}
\begin{longtable}{|p{2.2cm}|p{9.5cm}|p{2.3cm}|p{1.5cm}|}
    \caption{Test di sistema} \label{tab:ts} \\
    \hline
    \rowcolor{gray!20}
    \textbf{Codice} & \textbf{Descrizione} & \textbf{Riferimento} & \textbf{Stato} \\
    \hline
    \endfirsthead
    \hline
    \rowcolor{gray!20}
    \textbf{Codice} & \textbf{Descrizione} & \textbf{Riferimento} & \textbf{Stato} \\
    \hline
    \endhead
    \hline
    \endfoot
    \hline
    \endlastfoot
    \multicolumn{4}{|l|}{\cellcolor{gray!12}\textbf{Requisiti funzionali}} \\
    \hline
    TS-F-001 & Verifica che il sistema permetta al \glslink{redattore}{Redattore} di compilare il \glslink{prompt}{prompt} testuale descrittivo del contenuto da generare. & RF1-OB & NI \\
    \hline
    TS-F-002 & Verifica che il sistema permetta al \glslink{redattore}{Redattore} di selezionare il \glslink{tono-comunicativo}{tono comunicativo} della \glslink{bozza-draft}{bozza} da generare. & RF2-OB & NI \\
    \hline
    TS-F-003 & Verifica che il sistema permetta al \glslink{redattore}{Redattore} di selezionare lo \glslink{stile-editoriale}{stile editoriale} della \glslink{bozza-draft}{bozza} da generare. & RF3-OB & NI \\
    \hline
    TS-F-004 & Verifica che il sistema permetta al \glslink{redattore}{Redattore} di selezionare la lunghezza desiderata per la \glslink{bozza-draft}{bozza} da generare. & RF4-OP & NI \\
    \hline
    TS-F-005 & Verifica che il sistema permetta al \glslink{redattore}{Redattore} di salvare la configurazione corrente del \glslink{prompt}{prompt} nella libreria personale. & RF5-OB & NI \\
    \hline
    TS-F-006 & Verifica che il sistema permetta al \glslink{redattore}{Redattore} di svuotare i campi di configurazione ripristinando i valori predefiniti. & RF6-OB & NI \\
    \hline
    TS-F-007 & Verifica che il sistema avvii la generazione del contenuto tramite il modulo \glslink{ai-post-generator}{AI Post Generator} a partire dal \glslink{prompt}{prompt} e dai parametri configurati. & RF7-OB & NI \\
    \hline
    TS-F-008 & Verifica che il sistema notifichi il \glslink{redattore}{Redattore} qualora il campo \glslink{prompt}{prompt} sia assente o insufficiente prima dell'avvio della generazione. & RF8-OB & NI \\
    \hline
    TS-F-009 & Verifica che il sistema notifichi il \glslink{redattore}{Redattore} in caso di errore interno o \glslink{timeout}{timeout} durante la generazione iniziale del contenuto tramite \glslink{intelligenza-artificiale-ai}{AI}. & RF9-OB & NI \\
    \hline
    TS-F-010 & Verifica che il sistema notifichi il \glslink{redattore}{Redattore} in caso di errore interno o \glslink{timeout}{timeout} durante la rigenerazione del contenuto tramite \glslink{intelligenza-artificiale-ai}{AI}. & RF10-OB & NI \\
    \hline
    TS-F-011 & Verifica che il sistema visualizzi un'anteprima della \glslink{bozza-draft}{bozza} generata. & RF11-OB & NI \\
    \hline
    TS-F-012 & Verifica che il sistema visualizzi il titolo della \glslink{bozza-draft}{bozza} generata nell'anteprima. & RF12-OB & NI \\
    \hline
    TS-F-013 & Verifica che il sistema visualizzi il testo della \glslink{bozza-draft}{bozza} generata nell'anteprima. & RF13-OB & NI \\
    \hline
    TS-F-014 & Verifica che il sistema visualizzi l'immagine di copertina della \glslink{bozza-draft}{bozza} generata nell'anteprima. & RF14-OB & NI \\
    \hline
    TS-F-015 & Verifica che il sistema permetta al \glslink{redattore}{Redattore} di modificare manualmente il titolo della \glslink{bozza-draft}{bozza}. & RF15-OB & NI \\
    \hline
    TS-F-016 & Verifica che il sistema permetta al \glslink{redattore}{Redattore} di modificare manualmente il testo della \glslink{bozza-draft}{bozza}. & RF16-OB & NI \\
    \hline
    TS-F-017 & Verifica che il sistema notifichi il \glslink{redattore}{Redattore} in caso di formato file non valido durante la sostituzione dell'immagine di copertina. & RF17-OB & NI \\
    \hline
    TS-F-018 & Verifica che il sistema permetta al \glslink{redattore}{Redattore} di richiedere la rigenerazione della \glslink{bozza-draft}{bozza} mantenendo i parametri di configurazione. & RF18-OB & NI \\
    \hline
    TS-F-019 & Verifica che il sistema permetta al \glslink{redattore}{Redattore} di eliminare la \glslink{bozza-draft}{bozza} generata corrente. & RF19-OB & NI \\
    \hline
    TS-F-020 & Verifica che il sistema permetta al \glslink{redattore}{Redattore} di annullare le modifiche manuali apportate alla \glslink{bozza-draft}{bozza}, ripristinando il contenuto originale. & RF20-OP & NI \\
    \hline
    TS-F-021 & Verifica che il sistema notifichi il \glslink{redattore}{Redattore} qualora il contenuto della \glslink{bozza-draft}{bozza} risulti vuoto a seguito di modifiche. & RF21-OB & NI \\
    \hline
    TS-F-022 & Verifica che il sistema permetta al \glslink{redattore}{Redattore} di approvare e finalizzare il documento generato. & RF22-OB & NI \\
    \hline
    TS-F-023 & Verifica che il sistema visualizzi le opzioni disponibili per la finalizzazione del documento. & RF23-OB & NI \\
    \hline
    TS-F-024 & Verifica che il sistema permetta il salvataggio del documento come \glslink{bozza-draft}{bozza} nello storico. & RF24-OB & NI \\
    \hline
    TS-F-025 & Verifica che il sistema permetta l'esportazione del documento finalizzato in formato PDF. & RF25-OB & NI \\
    \hline
    TS-F-026 & Verifica che il sistema permetta al \glslink{redattore}{Redattore} di visualizzare un'anteprima del documento prima dell'esportazione definitiva. & RF26-OP & NI \\
    \hline
    TS-F-027 & Verifica che il sistema notifichi il \glslink{redattore}{Redattore} in caso di indirizzo e-mail non valido durante l'inoltro del documento. & RF27-OB & NI \\
    \hline
    TS-F-028 & Verifica che il sistema fornisca al \glslink{redattore}{Redattore} una libreria personale dei \glslink{prompt}{prompt} con il relativo storico delle generazioni passate. & RF28-OB & NI \\
    \hline
    TS-F-029 & Verifica che il sistema permetta al \glslink{redattore}{Redattore} di visualizzare l'elenco delle generazioni passate e i \glslink{prompt}{prompt} salvati nella libreria personale. & RF29-OB & NI \\
    \hline
    TS-F-030 & Verifica che il sistema permetta la ricerca e il filtraggio dello storico dei \glslink{prompt}{prompt} e delle generazioni passate. & RF30-OB & NI \\
    \hline
    TS-F-031 & Verifica che il sistema avvisi l'utente quando la ricerca nello storico non produce risultati. & RF31-OB & NI \\
    \hline
    TS-F-032 & Verifica che il sistema permetta il caricamento e il riutilizzo di una configurazione \glslink{prompt}{prompt} dalla libreria personale. & RF32-OB & NI \\
    \hline
    TS-F-033 & Verifica che il sistema permetta l'eliminazione di elementi dallo storico e dalla libreria dei \glslink{prompt}{prompt}. & RF33-OP & NI \\
    \hline
    TS-F-034 & Verifica che il sistema fornisca un modulo di valutazione qualitativa delle bozze generate. & RF34-OB & NI \\
    \hline
    TS-F-035 & Verifica che il sistema permetta al \glslink{redattore}{Redattore} di accedere al modulo di valutazione qualitativa della \glslink{bozza-draft}{bozza} generata. & RF35-OB & NI \\
    \hline
    TS-F-036 & Verifica che il sistema permetta al \glslink{redattore}{Redattore} di assegnare un \glslink{rating}{rating} numerico alla \glslink{bozza-draft}{bozza} generata. & RF36-OB & NI \\
    \hline
    TS-F-037 & Verifica che il sistema permetta al \glslink{redattore}{Redattore} di inviare la valutazione al sistema. & RF37-OB & NI \\
    \hline
    TS-F-038 & Verifica che il sistema notifichi il \glslink{redattore}{Redattore} qualora tenti di inviare la valutazione senza aver assegnato un \glslink{rating}{rating}. & RF38-OB & NI \\
    \hline
    TS-F-039 & Verifica che il sistema notifichi il \glslink{redattore}{Redattore} qualora il commento testuale superi il limite di caratteri consentito. & RF39-OB & NI \\
    \hline
    TS-F-040 & Verifica che il sistema fornisca all'\glslink{analista-dati}{Analista dati} una \glslink{dashboard}{dashboard} dedicata con metriche di utilizzo del modulo \glslink{ai-assistant-generativo}{AI Assistant Generativo}. & RF40-OB & NI \\
    \hline
    TS-F-041 & Verifica che il sistema visualizzi nella \glslink{dashboard}{dashboard} una panoramica dell'attività complessiva del modulo \glslink{intelligenza-artificiale-ai}{AI} Assistant. & RF41-OB & NI \\
    \hline
    TS-F-042 & Verifica che il sistema visualizzi nella \glslink{dashboard}{dashboard} le statistiche di utilizzo del modulo \glslink{intelligenza-artificiale-ai}{AI} Assistant. & RF42-OB & NI \\
    \hline
    TS-F-043 & Verifica che il sistema visualizzi nella \glslink{dashboard}{dashboard} il \glslink{rating}{rating} medio delle bozze generate. & RF43-OB & NI \\
    \hline
    TS-F-044 & Verifica che il sistema permetta all'\glslink{analista-dati}{Analista dati} di consultare i feedback testuali ricevuti dagli utenti. & RF44-OB & NI \\
    \hline
    TS-F-045 & Verifica che il sistema notifichi l'\glslink{analista-dati}{Analista dati} qualora i dati non siano disponibili per l'esportazione del report. & RF45-OB & NI \\
    \hline
    TS-F-046 & Verifica che il sistema permetta al \glslink{consulente-del-lavoro-cdl}{Consulente del lavoro} di caricare singoli documenti o lotti nel \glslink{repository}{repository} dedicato. & RF46-OB & NI \\
    \hline
    TS-F-047 & Verifica che il sistema rilevi i documenti duplicati durante l'upload e notifichi il \glslink{consulente-del-lavoro-cdl}{Consulente del lavoro} richiedendone conferma. & RF47-OB & NI \\
    \hline
    TS-F-048 & Verifica che il sistema rilevi i formati file non supportati durante l'upload e notifichi il \glslink{consulente-del-lavoro-cdl}{Consulente del lavoro}. & RF48-OB & NI \\
    \hline
    TS-F-049 & Verifica che il sistema classifichi automaticamente i documenti caricati e popoli i \glslink{metadati}{metadati} associati tramite il modulo \glslink{ai-analyst}{AI Analyst}. & RF49-OB & NI \\
    \hline
    TS-F-050 & Verifica che il sistema permetta al \glslink{consulente-del-lavoro-cdl}{Consulente del lavoro} di visualizzare lo storico dei documenti analizzati. & RF50-OB & NI \\
    \hline
    TS-F-051 & Verifica che il sistema notifichi l'utente quando nessun documento corrisponde \glslink{intelligenza-artificiale-ai}{ai} criteri di visualizzazione dello storico. & RF51-OB & NI \\
    \hline
    TS-F-052 & Verifica che il sistema permetta il filtraggio dello storico documenti tramite criteri multipli. & RF52-OB & NI \\
    \hline
    TS-F-053 & Verifica che il sistema permetta il filtraggio dello storico documenti tramite barra di ricerca testuale. & RF53-OB & NI \\
    \hline
    TS-F-054 & Verifica che il sistema permetta il filtraggio dello storico documenti per stato di invio. & RF54-OB & NI \\
    \hline
    TS-F-055 & Verifica che il sistema permetta il filtraggio dello storico documenti per livello di \glslink{confidenza-confidence-score}{confidenza} dell'analisi \glslink{intelligenza-artificiale-ai}{AI}. & RF55-OB & NI \\
    \hline
    TS-F-056 & Verifica che il sistema permetta al \glslink{consulente-del-lavoro-cdl}{Consulente del lavoro} di visualizzare i dettagli di un singolo documento dallo storico. & RF56-OB & NI \\
    \hline
    TS-F-057 & Verifica che il sistema permetta la visualizzazione del documento originale nella versione splittata. & RF57-OB & NI \\
    \hline
    TS-F-058 & Verifica che il sistema permetta la visualizzazione del documento originale non splittato. & RF58-OB & NI \\
    \hline
    TS-F-059 & Verifica che il sistema visualizzi il nome e cognome del destinatario associato al documento. & RF59-OB & NI \\
    \hline
    TS-F-060 & Verifica che il sistema visualizzi l'azienda associata al documento. & RF60-OB & NI \\
    \hline
    TS-F-061 & Verifica che il sistema visualizzi il nome del file del documento. & RF61-OB & NI \\
    \hline
    TS-F-062 & Verifica che il sistema visualizzi la data del documento. & RF62-OB & NI \\
    \hline
    TS-F-063 & Verifica che il sistema visualizzi il numero di pagine del documento. & RF63-OB & NI \\
    \hline
    TS-F-064 & Verifica che il sistema visualizzi la tipologia del documento. & RF64-OB & NI \\
    \hline
    TS-F-065 & Verifica che il sistema visualizzi la breve descrizione del documento. & RF65-OB & NI \\
    \hline
    TS-F-066 & Verifica che il sistema visualizzi il livello di \glslink{confidenza-confidence-score}{confidenza} in percentuale del documento. & RF66-OB & NI \\
    \hline
    TS-F-067 & Verifica che il sistema visualizzi lo stato di invio del documento. & RF67-OB & NI \\
    \hline
    TS-F-068 & Verifica che il sistema visualizzi l'indirizzo e-mail del destinatario del documento. & RF68-OP & NI \\
    \hline
    TS-F-069 & Verifica che il sistema visualizzi il codice fiscale del destinatario. & RF69-OP & NI \\
    \hline
    TS-F-070 & Verifica che il sistema visualizzi la matricola del dipendente. & RF70-OP & NI \\
    \hline
    TS-F-071 & Verifica che il sistema visualizzi la data e ora di caricamento del documento. & RF71-OP & NI \\
    \hline
    TS-F-072 & Verifica che il sistema gestisca l'assenza del nome del destinatario visualizzando un indicatore apposito. & RF72-OB & NI \\
    \hline
    TS-F-073 & Verifica che il sistema gestisca l'assenza dell'azienda associata al documento visualizzando un indicatore apposito. & RF73-OB & NI \\
    \hline
    TS-F-074 & Verifica che il sistema gestisca l'assenza della data del documento visualizzando un indicatore apposito. & RF74-OB & NI \\
    \hline
    TS-F-075 & Verifica che il sistema gestisca l'assenza della tipologia del documento visualizzando un indicatore apposito. & RF75-OB & NI \\
    \hline
    TS-F-076 & Verifica che il sistema gestisca l'assenza della breve descrizione visualizzando un indicatore apposito. & RF76-OB & NI \\
    \hline
    TS-F-077 & Verifica che il sistema gestisca l'assenza del livello di \glslink{confidenza-confidence-score}{confidenza} visualizzando un indicatore apposito. & RF77-OB & NI \\
    \hline
    TS-F-078 & Verifica che il sistema permetta al \glslink{consulente-del-lavoro-cdl}{Consulente del lavoro} di modificare i \glslink{metadati}{metadati} di un documento presente nello storico. & RF78-OB & NI \\
    \hline
    TS-F-079 & Verifica che il sistema permetta la modifica del nome e cognome del destinatario associato al documento. & RF79-OB & NI \\
    \hline
    TS-F-080 & Verifica che il sistema permetta la modifica dell'azienda associata al documento. & RF80-OB & NI \\
    \hline
    TS-F-081 & Verifica che il sistema permetta la modifica della data del documento. & RF81-OB & NI \\
    \hline
    TS-F-082 & Verifica che il sistema permetta la modifica della tipologia del documento. & RF82-OB & NI \\
    \hline
    TS-F-083 & Verifica che il sistema permetta la modifica della breve descrizione del documento. & RF83-OB & NI \\
    \hline
    TS-F-084 & Verifica che il sistema permetta la modifica dell'indirizzo e-mail del destinatario. & RF84-OP & NI \\
    \hline
    TS-F-085 & Verifica che il sistema permetta la modifica del codice fiscale del destinatario. & RF85-OP & NI \\
    \hline
    TS-F-086 & Verifica che il sistema permetta la modifica della matricola del dipendente. & RF86-OP & NI \\
    \hline
    TS-F-087 & Verifica che il sistema notifichi il \glslink{consulente-del-lavoro-cdl}{Consulente del lavoro} in caso di formato e-mail non valido durante la modifica dei \glslink{metadati}{metadati}. & RF87-OP & NI \\
    \hline
    TS-F-088 & Verifica che il sistema notifichi il \glslink{consulente-del-lavoro-cdl}{Consulente del lavoro} in caso di codice fiscale non valido durante la modifica dei \glslink{metadati}{metadati}. & RF88-OP & NI \\
    \hline
    TS-F-089 & Verifica che il sistema generi automaticamente il contenuto del messaggio di invio per il documento da spedire. & RF89-OB & NI \\
    \hline
    TS-F-090 & Verifica che il sistema visualizzi il destinatario del messaggio di invio generato automaticamente. & RF90-OB & NI \\
    \hline
    TS-F-091 & Verifica che il sistema visualizzi l'oggetto del messaggio di invio generato automaticamente. & RF91-OB & NI \\
    \hline
    TS-F-092 & Verifica che il sistema visualizzi il testo del messaggio di invio generato automaticamente. & RF92-OB & NI \\
    \hline
    TS-F-093 & Verifica che il sistema permetta la conferma e l'effettivo invio del messaggio al destinatario. & RF93-OB & NI \\
    \hline
    TS-F-094 & Verifica che il sistema permetta al \glslink{consulente-del-lavoro-cdl}{Consulente del lavoro} di modificare il contenuto del messaggio di invio generato prima della spedizione. & RF94-OB & NI \\
    \hline
    TS-F-095 & Verifica che il sistema permetta la modifica del destinatario del messaggio di invio. & RF95-OB & NI \\
    \hline
    TS-F-096 & Verifica che il sistema permetta la modifica dell'oggetto del messaggio di invio. & RF96-OB & NI \\
    \hline
    TS-F-097 & Verifica che il sistema permetta la modifica del testo del messaggio di invio. & RF97-OB & NI \\
    \hline
    TS-F-098 & Verifica che il sistema fornisca all'\glslink{analista-dati}{Analista dati} una \glslink{dashboard}{dashboard} dedicata con metriche di monitoraggio del modulo \glslink{intelligenza-artificiale-ai}{AI} Co-Pilot. & RF98-OB & NI \\
    \hline
    TS-F-099 & Verifica che il sistema visualizzi nella \glslink{dashboard}{dashboard} del Co-Pilot il numero totale di documenti analizzati. & RF99-OB & NI \\
    \hline
    TS-F-100 & Verifica che il sistema visualizzi nella \glslink{dashboard}{dashboard} del Co-Pilot la percentuale di classificazioni corrette. & RF100-OB & NI \\
    \hline
    TS-F-101 & Verifica che il sistema visualizzi nella \glslink{dashboard}{dashboard} del Co-Pilot i documenti sotto \glslink{soglia-di-confidenza}{soglia di confidenza}. & RF101-OB & NI \\
    \hline
    TS-F-102 & Verifica che il sistema visualizzi nella \glslink{dashboard}{dashboard} del Co-Pilot la percentuale di destinatari riconosciuti automaticamente. & RF102-OB & NI \\
    \hline
    TS-F-103 & Verifica che il sistema visualizzi nella \glslink{dashboard}{dashboard} del Co-Pilot il tempo medio di analisi dei documenti. & RF103-OB & NI \\
    \hline
    \multicolumn{4}{|l|}{\cellcolor{gray!12}\textbf{Requisiti di qualità}} \\
    \hline
    TS-Q-001 & Verifica che l'\glslink{analisi-dei-requisiti-adr}{Analisi dei Requisiti} sia completa di casi d'uso modellati tramite UML. & RQ1-OB & NI \\
    \hline
    TS-Q-002 & Verifica che le pratiche di sviluppo siano conformi alle \glslink{norme-di-progetto}{Norme di Progetto} adottate dal gruppo. & RQ2-OB & NI \\
    \hline
    TS-Q-003 & Verifica che il prodotto superi tutti i test automatizzati con una copertura del codice non inferiore all'80\%. & RQ3-OB & NI \\
    \hline
    TS-Q-004 & Verifica che la documentazione del codice sorgente rispetti le convenzioni definite nelle \glslink{norme-di-progetto}{Norme di Progetto}. & RQ4-OB & NI \\
    \hline
    TS-Q-005 & Verifica che il sorgente sia gestito tramite un sistema di versionamento con istruzioni di setup incluse nel \glslink{repository}{repository}. & RQ5-OB & NI \\
    \hline
    TS-Q-006 & Verifica che il manuale utente del prodotto sia prodotto e aggiornato. & RQ6-OB & NI \\
    \hline
    TS-Q-007 & Verifica tramite sessione di user testing che gli utenti siano in grado di apprendere le funzionalità principali del sistema in meno di 30 minuti. & MPD11 & NI \\
    \hline
    \multicolumn{4}{|l|}{\cellcolor{gray!12}\textbf{Requisiti di vincolo}} \\
    \hline
    TS-V-001 & Verifica che il codice sorgente sia versionato tramite \glslink{git}{Git} su un \glslink{repository}{repository} condiviso. & RVC1-OB & NI \\
    \hline
    TS-V-002 & Verifica che il codice sorgente rispetti le regole di formattazione definite da \glslink{laravel}{Laravel} Pint. & RVC2-OB & NI \\
    \hline
    TS-V-003 & Verifica che la suite di test automatizzati sia scritta con il framework \glslink{pest}{Pest}. & RVC3-OB & NI \\
    \hline
    TS-V-004 & Verifica che il backend sia sviluppato con il framework \glslink{laravel}{Laravel} 12 su PHP 8.4. & RVC4-OB & NI \\
    \hline
    TS-V-005 & Verifica che il database relazionale del sistema sia \glslink{postgresql}{PostgreSQL}. & RVC5-OB & NI \\
    \hline
    TS-V-006 & Verifica che le code asincrone e la gestione dei job in background utilizzino \glslink{redis}{Redis}. & RVC6-OB & NI \\
    \hline
    TS-V-007 & Verifica che il frontend delle interfacce utente sia realizzato tramite il motore di template Blade con CSS e JavaScript vanilla. & RVC7-OB & NI \\
    \hline
    TS-V-008 & Verifica che l'ambiente di sviluppo e di esecuzione del sistema sia containerizzato tramite Docker e orchestrato con Docker Compose. & RVC8-OB & NI \\
    \hline
    TS-V-009 & Verifica che il web server dell'applicazione sia Nginx. & RVC9-OB & NI \\
    \hline
    TS-V-010 & Verifica che lo storage dei file caricati nel sistema utilizzi \glslink{minio}{MinIO}, servizio compatibile con le API Amazon S3. & RVC10-OB & NI \\
    \hline
    TS-V-011 & Verifica che i modelli di \glslink{intelligenza-artificiale-ai}{intelligenza artificiale} generativa possano essere erogati tramite \glslink{amazon-bedrock}{Amazon Bedrock}. & RVC11-OP & NI \\
    \hline
    TS-V-012 & Verifica che il sistema superi i controlli di accessibilità automatizzati tramite pa11y e axe sulle interfacce principali. & RVC12-OB & NI \\
    \hline
    TS-V-013 & Verifica che il sistema sia fruibile sui principali browser moderni evergreen nelle ultime due versioni stabili. & RVC13-OB & NI \\
    \hline
    TS-V-014 & Verifica che l'ambiente di esecuzione sia containerizzato con Docker Compose con lo stack completo richiesto. & RVC14-OB & NI \\
    \hline
    TS-V-015 & Verifica che la configurazione minima del server rispetti i requisiti di 2 vCPU, 4 GB di RAM e 10 GB di spazio disco. & RVC15-OB & NI \\
    \hline
    TS-V-016 & Verifica che l'utilizzo delle funzionalità \glslink{intelligenza-artificiale-ai}{AI} richieda credenziali AWS valide e accesso outbound ad \glslink{amazon-bedrock}{Amazon Bedrock}. & RVC16-OB & NI \\
    \hline
\end{longtable}

\subsection{Test di accettazione}
I \glslink{test-di-accettazione}{test di accettazione} verificano che il sistema soddisfi i requisiti utente relativi al \glslink{capitolato}{capitolato} d'appalto. Ogni test è associato al \glslink{caso-d-uso-uc}{caso d'uso} di riferimento.

\renewcommand{\arraystretch}{1.15}
\rowcolors{2}{gray!8}{white}
\begin{longtable}{|p{2.2cm}|p{9.5cm}|p{2.3cm}|p{1.5cm}|}
    \caption{Test di accettazione} \label{tab:ta} \\
    \hline
    \rowcolor{gray!20}
    \textbf{Codice} & \textbf{Descrizione} & \textbf{Riferimento} & \textbf{Stato} \\
    \hline
    \endfirsthead
    \hline
    \rowcolor{gray!20}
    \textbf{Codice} & \textbf{Descrizione} & \textbf{Riferimento} & \textbf{Stato} \\
    \hline
    \endhead
    \hline
    \endfoot
    \hline
    \endlastfoot
    TA-001 & Verifica che il sistema permetta al \glslink{redattore}{Redattore} di configurare completamente i parametri del \glslink{prompt}{prompt} (tono, stile, lunghezza) e avviare la generazione \glslink{intelligenza-artificiale-ai}{AI} con successo. & UC-1 & NI \\
    \hline
    TA-002 & Verifica che il sistema notifichi il \glslink{redattore}{Redattore} qualora tenti di avviare la generazione con il campo \glslink{prompt}{prompt} assente o insufficiente. & UC-1 & NI \\
    \hline
    TA-003 & Verifica che il sistema visualizzi correttamente l'anteprima della \glslink{bozza-draft}{bozza} generata con titolo, testo e immagine di copertina. & UC-2 & NI \\
    \hline
    TA-004 & Verifica che il sistema permetta al \glslink{redattore}{Redattore} di modificare manualmente il titolo e il testo della \glslink{bozza-draft}{bozza} e di visualizzare il contenuto aggiornato. & UC-2 & NI \\
    \hline
    TA-005 & Verifica che il sistema permetta al \glslink{redattore}{Redattore} di rigenerare la \glslink{bozza-draft}{bozza} mantenendo la configurazione del \glslink{prompt}{prompt} corrente. & UC-2 & NI \\
    \hline
    TA-006 & Verifica che il sistema permetta al \glslink{redattore}{Redattore} di eliminare la \glslink{bozza-draft}{bozza} generata corrente. & UC-2 & NI \\
    \hline
    TA-007 & Verifica che il sistema permetta il salvataggio del documento finalizzato come \glslink{bozza-draft}{bozza} nello storico. & UC-3 & NI \\
    \hline
    TA-008 & Verifica che il sistema permetta l'esportazione del documento finalizzato in formato PDF, con anteprima pre-esportazione. & UC-3 & NI \\
    \hline
    TA-009 & Verifica che il sistema permetta l'inoltro del documento finalizzato via e-mail e notifichi il \glslink{redattore}{Redattore} in caso di indirizzo non valido. & UC-3 & NI \\
    \hline
    TA-010 & Verifica che il sistema visualizzi la libreria personale con lo storico delle generazioni e i \glslink{prompt}{prompt} salvati, con possibilità di eliminazione. & UC-4 & NI \\
    \hline
    TA-011 & Verifica che il sistema permetta la ricerca, il filtraggio e il riutilizzo di una configurazione \glslink{prompt}{prompt} dalla libreria personale. & UC-4 & NI \\
    \hline
    TA-012 & Verifica che il sistema permetta al \glslink{redattore}{Redattore} di inserire e inviare una valutazione qualitativa con \glslink{rating}{rating} numerico sulla \glslink{bozza-draft}{bozza} generata. & UC-5 & NI \\
    \hline
    TA-013 & Verifica che il sistema notifichi il \glslink{redattore}{Redattore} qualora tenti di inviare la valutazione senza aver assegnato un \glslink{rating}{rating}. & UC-5 & NI \\
    \hline
    TA-014 & Verifica che il sistema esponga all'\glslink{analista-dati}{Analista dati} la \glslink{dashboard}{dashboard} \glslink{intelligenza-artificiale-ai}{AI} Assistant con panoramica dell'attività complessiva e statistiche di utilizzo. & UC-6 & NI \\
    \hline
    TA-015 & Verifica che il sistema visualizzi nella \glslink{dashboard}{dashboard} \glslink{intelligenza-artificiale-ai}{AI} Assistant il \glslink{rating}{rating} medio delle bozze e permetta la consultazione dei feedback testuali. & UC-6 & NI \\
    \hline
    TA-016 & Verifica che il sistema permetta al \glslink{consulente-del-lavoro-cdl}{Consulente del lavoro} di caricare un singolo documento nel \glslink{repository}{repository} con classificazione automatica tramite \glslink{ai-analyst}{AI Analyst}. & UC-7, UC-8 & NI \\
    \hline
    TA-017 & Verifica che il sistema permetta l'upload di un lotto di documenti e rilevi i duplicati richiedendo conferma al \glslink{consulente-del-lavoro-cdl}{Consulente del lavoro}. & UC-7 & NI \\
    \hline
    TA-018 & Verifica che il sistema rilevi un file in formato non supportato durante l'upload e notifichi il \glslink{consulente-del-lavoro-cdl}{Consulente del lavoro}. & UC-7 & NI \\
    \hline
    TA-019 & Verifica che il sistema visualizzi al \glslink{consulente-del-lavoro-cdl}{Consulente del lavoro} lo storico dei documenti analizzati e notifichi in caso di assenza di risultati. & UC-9, UC-9.E1 & NI \\
    \hline
    TA-020 & Verifica che il sistema permetta il filtraggio dello storico documenti tramite la combinazione di più criteri contemporaneamente. & UC-10 & NI \\
    \hline
    TA-021 & Verifica che il sistema permetta il filtraggio dello storico documenti tramite la barra di ricerca testuale. & UC-11 & NI \\
    \hline
    TA-022 & Verifica che il sistema permetta il filtraggio dello storico documenti per stato di invio. & UC-12 & NI \\
    \hline
    TA-023 & Verifica che il sistema permetta il filtraggio dello storico documenti per livello di \glslink{confidenza-confidence-score}{confidenza} dell'analisi \glslink{intelligenza-artificiale-ai}{AI}. & UC-13 & NI \\
    \hline
    TA-024 & Verifica che il sistema permetta il filtraggio dello storico documenti per mese e anno. & UC-14 & NI \\
    \hline
    TA-025 & Verifica che il sistema visualizzi i dettagli completi di un documento dallo storico, inclusi i \glslink{metadati}{metadati} associati e la versione splittata del documento originale. & UC-15 & NI \\
    \hline
    TA-026 & Verifica che il sistema permetta la modifica dei \glslink{metadati}{metadati} principali di un documento (nome destinatario, azienda, data, tipologia, descrizione) e il relativo salvataggio. & UC-16 & NI \\
    \hline
    TA-027 & Verifica che il sistema notifichi il \glslink{consulente-del-lavoro-cdl}{Consulente del lavoro} in caso di e-mail o codice fiscale in formato non valido durante la modifica dei \glslink{metadati}{metadati}. & UC-16 & NI \\
    \hline
    TA-028 & Verifica che il sistema generi automaticamente il messaggio di invio per il documento selezionato e ne visualizzi destinatario, oggetto e testo. & UC-17 & NI \\
    \hline
    TA-029 & Verifica che il sistema permetta la modifica del contenuto del messaggio di invio generato e la conferma della spedizione al destinatario. & UC-18 & NI \\
    \hline
    TA-030 & Verifica che il sistema esponga all'\glslink{analista-dati}{Analista dati} la \glslink{dashboard}{dashboard} Co-Pilot con metriche di monitoraggio dell'analisi documentale (totale documenti, classificazioni corrette, documenti sotto soglia, destinatari riconosciuti, tempo medio di analisi). & UC-19 & NI \\
    \hline
\end{longtable}

\subsection{Tracciamento}
La tabella seguente riporta la corrispondenza tra i \glslink{test-di-sistema}{test di sistema} e i requisiti coperti, e tra i \glslink{test-di-accettazione}{test di accettazione} e i casi d'uso di riferimento.

{\footnotesize
\renewcommand{\arraystretch}{1.1}
\rowcolors{2}{gray!8}{white}
\begin{longtable}{|p{2.5cm}|p{2.2cm}|p{2.5cm}|p{2.2cm}|}
    \caption{Tracciamento test di sistema -- requisiti} \label{tab:tracciamento-ts} \\
    \hline
    \rowcolor{gray!20}
    \textbf{Codice test} & \textbf{Requisito} & \textbf{Codice test} & \textbf{Requisito} \\
    \hline
    \endfirsthead
    \hline
    \rowcolor{gray!20}
    \textbf{Codice test} & \textbf{Requisito} & \textbf{Codice test} & \textbf{Requisito} \\
    \hline
    \endhead
    \hline
    \endfoot
    \hline
    \endlastfoot
    \multicolumn{4}{|l|}{\cellcolor{gray!12}\textbf{Funzionali}} \\
    \hline
    TS-F-001 & RF1-OB & TS-F-002 & RF2-OB \\
    \hline
    TS-F-003 & RF3-OB & TS-F-004 & RF4-OP \\
    \hline
    TS-F-005 & RF5-OB & TS-F-006 & RF6-OB \\
    \hline
    TS-F-007 & RF7-OB & TS-F-008 & RF8-OB \\
    \hline
    TS-F-009 & RF9-OB & TS-F-010 & RF10-OB \\
    \hline
    TS-F-011 & RF11-OB & TS-F-012 & RF12-OB \\
    \hline
    TS-F-013 & RF13-OB & TS-F-014 & RF14-OB \\
    \hline
    TS-F-015 & RF15-OB & TS-F-016 & RF16-OB \\
    \hline
    TS-F-017 & RF17-OB & TS-F-018 & RF18-OB \\
    \hline
    TS-F-019 & RF19-OB & TS-F-020 & RF20-OP \\
    \hline
    TS-F-021 & RF21-OB & TS-F-022 & RF22-OB \\
    \hline
    TS-F-023 & RF23-OB & TS-F-024 & RF24-OB \\
    \hline
    TS-F-025 & RF25-OB & TS-F-026 & RF26-OP \\
    \hline
    TS-F-027 & RF27-OB & TS-F-028 & RF28-OB \\
    \hline
    TS-F-029 & RF29-OB & TS-F-030 & RF30-OB \\
    \hline
    TS-F-031 & RF31-OB & TS-F-032 & RF32-OB \\
    \hline
    TS-F-033 & RF33-OP & TS-F-034 & RF34-OB \\
    \hline
    TS-F-035 & RF35-OB & TS-F-036 & RF36-OB \\
    \hline
    TS-F-037 & RF37-OB & TS-F-038 & RF38-OB \\
    \hline
    TS-F-039 & RF39-OB & TS-F-040 & RF40-OB \\
    \hline
    TS-F-041 & RF41-OB & TS-F-042 & RF42-OB \\
    \hline
    TS-F-043 & RF43-OB & TS-F-044 & RF44-OB \\
    \hline
    TS-F-045 & RF45-OB & TS-F-046 & RF46-OB \\
    \hline
    TS-F-047 & RF47-OB & TS-F-048 & RF48-OB \\
    \hline
    TS-F-049 & RF49-OB & TS-F-050 & RF50-OB \\
    \hline
    TS-F-051 & RF51-OB & TS-F-052 & RF52-OB \\
    \hline
    TS-F-053 & RF53-OB & TS-F-054 & RF54-OB \\
    \hline
    TS-F-055 & RF55-OB & TS-F-056 & RF56-OB \\
    \hline
    TS-F-057 & RF57-OB & TS-F-058 & RF58-OB \\
    \hline
    TS-F-059 & RF59-OB & TS-F-060 & RF60-OB \\
    \hline
    TS-F-061 & RF61-OB & TS-F-062 & RF62-OB \\
    \hline
    TS-F-063 & RF63-OB & TS-F-064 & RF64-OB \\
    \hline
    TS-F-065 & RF65-OB & TS-F-066 & RF66-OB \\
    \hline
    TS-F-067 & RF67-OB & TS-F-068 & RF68-OP \\
    \hline
    TS-F-069 & RF69-OP & TS-F-070 & RF70-OP \\
    \hline
    TS-F-071 & RF71-OP & TS-F-072 & RF72-OB \\
    \hline
    TS-F-073 & RF73-OB & TS-F-074 & RF74-OB \\
    \hline
    TS-F-075 & RF75-OB & TS-F-076 & RF76-OB \\
    \hline
    TS-F-077 & RF77-OB & TS-F-078 & RF78-OB \\
    \hline
    TS-F-079 & RF79-OB & TS-F-080 & RF80-OB \\
    \hline
    TS-F-081 & RF81-OB & TS-F-082 & RF82-OB \\
    \hline
    TS-F-083 & RF83-OB & TS-F-084 & RF84-OP \\
    \hline
    TS-F-085 & RF85-OP & TS-F-086 & RF86-OP \\
    \hline
    TS-F-087 & RF87-OP & TS-F-088 & RF88-OP \\
    \hline
    TS-F-089 & RF89-OB & TS-F-090 & RF90-OB \\
    \hline
    TS-F-091 & RF91-OB & TS-F-092 & RF92-OB \\
    \hline
    TS-F-093 & RF93-OB & TS-F-094 & RF94-OB \\
    \hline
    TS-F-095 & RF95-OB & TS-F-096 & RF96-OB \\
    \hline
    TS-F-097 & RF97-OB & TS-F-098 & RF98-OB \\
    \hline
    TS-F-099 & RF99-OB & TS-F-100 & RF100-OB \\
    \hline
    TS-F-101 & RF101-OB & TS-F-102 & RF102-OB \\
    \hline
    TS-F-103 & RF103-OB & & \\
    \hline
    \multicolumn{4}{|l|}{\cellcolor{gray!12}\textbf{Qualità}} \\
    \hline
    TS-Q-001 & RQ1-OB & TS-Q-002 & RQ2-OB \\
    \hline
    TS-Q-003 & RQ3-OB & TS-Q-004 & RQ4-OB \\
    \hline
    TS-Q-005 & RQ5-OB & TS-Q-006 & RQ6-OB \\
    \hline
    TS-Q-007 & MPD11 & & \\
    \hline
    \multicolumn{4}{|l|}{\cellcolor{gray!12}\textbf{Vincolo}} \\
    \hline
    TS-V-001 & RVC1-OB & TS-V-002 & RVC2-OB \\
    \hline
    TS-V-003 & RVC3-OB & TS-V-004 & RVC4-OB \\
    \hline
    TS-V-005 & RVC5-OB & TS-V-006 & RVC6-OB \\
    \hline
    TS-V-007 & RVC7-OB & TS-V-008 & RVC8-OB \\
    \hline
    TS-V-009 & RVC9-OB & TS-V-010 & RVC10-OB \\
    \hline
    TS-V-011 & RVC11-OP & TS-V-012 & RVC12-OB \\
    \hline
    TS-V-013 & RVC13-OB & TS-V-014 & RVC14-OB \\
    \hline
    TS-V-015 & RVC15-OB & TS-V-016 & RVC16-OB \\
    \hline
\end{longtable}
}

{\footnotesize
\renewcommand{\arraystretch}{1.1}
\rowcolors{2}{gray!8}{white}
\begin{longtable}{|p{2.5cm}|p{2.2cm}|p{2.5cm}|p{2.2cm}|}
    \caption{Tracciamento test di accettazione -- casi d'uso} \label{tab:tracciamento-ta} \\
    \hline
    \rowcolor{gray!20}
    \textbf{Codice test} & \textbf{\glslink{caso-d-uso-uc}{Caso d'uso}} & \textbf{Codice test} & \textbf{\glslink{caso-d-uso-uc}{Caso d'uso}} \\
    \hline
    \endfirsthead
    \hline
    \rowcolor{gray!20}
    \textbf{Codice test} & \textbf{\glslink{caso-d-uso-uc}{Caso d'uso}} & \textbf{Codice test} & \textbf{\glslink{caso-d-uso-uc}{Caso d'uso}} \\
    \hline
    \endhead
    \hline
    \endfoot
    \hline
    \endlastfoot
    TA-001 & UC-1 & TA-002 & UC-1 \\
    \hline
    TA-003 & UC-2 & TA-004 & UC-2 \\
    \hline
    TA-005 & UC-2 & TA-006 & UC-2 \\
    \hline
    TA-007 & UC-3 & TA-008 & UC-3 \\
    \hline
    TA-009 & UC-3 & TA-010 & UC-4 \\
    \hline
    TA-011 & UC-4 & TA-012 & UC-5 \\
    \hline
    TA-013 & UC-5 & TA-014 & UC-6 \\
    \hline
    TA-015 & UC-6 & TA-016 & UC-7, UC-8 \\
    \hline
    TA-017 & UC-7 & TA-018 & UC-7 \\
    \hline
    TA-019 & UC-9, UC-9.E1 & TA-020 & UC-10 \\
    \hline
    TA-021 & UC-11 & TA-022 & UC-12 \\
    \hline
    TA-023 & UC-13 & TA-024 & UC-14 \\
    \hline
    TA-025 & UC-15 & TA-026 & UC-16 \\
    \hline
    TA-027 & UC-16 & TA-028 & UC-17 \\
    \hline
    TA-029 & UC-18 & TA-030 & UC-19 \\
    \hline
\end{longtable}
}


\section{Cruscotto di qualità}

Il cruscotto di qualità raccoglie i dati misurati nelle diverse iterazioni del progetto, mettendoli in relazione con le soglie accettabili e ottimali definite nella sezione \emph{Qualità di processo}. L'obiettivo è fornire una visione sintetica e aggiornata dello stato qualitativo del prodotto e dei processi, permettendo di identificare tempestivamente eventuali aree di miglioramento.

Le misurazioni si riferiscono \glslink{intelligenza-artificiale-ai}{ai} quattro \glslink{sprint}{sprint} del periodo \glslink{rtb-requirements-and-technology-baseline}{RTB}:

\begin{itemize}
    \item \textbf{\glslink{sprint}{Sprint} 1 (S1):} 30/03/2026 – 12/04/2026
    \item \textbf{\glslink{sprint}{Sprint} 2 (S2):} 13/04/2026 – 26/04/2026
    \item \textbf{\glslink{sprint}{Sprint} 3 (S3):} 27/04/2026 – 10/05/2026
    \item \textbf{\glslink{sprint}{Sprint} 4 (S4):} 11/05/2026 – 24/05/2026
\end{itemize}

Il Budget At Completion (BAC) dell'intero progetto è pari a \EUR{12.575,00}.

% ─────────────────────────────────────────────────────────────────────────────
\subsection{Processi primari — Fornitura}

\subsubsection{MPC-EV e MPC-PV — Earned Value e Planned Value}

Il Planned Value (PV) rappresenta il valore pianificato del lavoro da completare entro ciascuno \glslink{sprint}{sprint}, estratto dai preventivi del \glslink{piano-di-progetto}{Piano di Progetto}. L'Earned Value (EV) rappresenta il valore del lavoro effettivamente completato. Nella fase \glslink{rtb-requirements-and-technology-baseline}{RTB} si assume che gli obiettivi di ogni \glslink{sprint}{sprint} siano stati completamente raggiunti (come risulta dalle retrospettive nel \glslink{piano-di-progetto}{Piano di Progetto}), pertanto $\text{EV} = \text{PV}$ per tutti gli \glslink{sprint}{sprint}.

\begin{table}[H]
\centering
\rowcolors{2}{gray!15}{white}
\begin{tabular}{|l|r|r|r|r|}
    \hline
    \rowcolor{gray!35}
    \textbf{\glslink{sprint}{Sprint}} & \textbf{PV (\euro)} & \textbf{EV (\euro)} & \textbf{PV cum.\ (\euro)} & \textbf{EV cum.\ (\euro)} \\
    \hline
    S1 & 670,00  & 670,00  & 670,00  & 670,00  \\
    S2 & 565,00  & 565,00  & 1.235,00 & 1.235,00 \\
    S3 & 750,00  & 750,00  & 1.985,00 & 1.985,00 \\
    S4 & 850,00  & 850,00  & 2.835,00 & 2.835,00 \\
    \hline
\end{tabular}
\caption{Andamento MPC-EV e MPC-PV per sprint (fase RTB)}
\end{table}

\begin{figure}[H]
\centering
\begin{tikzpicture}
\begin{axis}[
    title={Andamento MPC-EV e MPC-PV Cumulativi},
    xlabel={\glslink{sprint}{Sprint}},
    ylabel={Valore (\euro)},
    symbolic x coords={S1, S2, S3, S4},
    xtick=data,
    ymin=0, ymax=3200,
    grid=both,
    grid style={dashed, gray!30},
    legend pos=north west,
    width=0.85\textwidth,
    height=6.5cm
]
\addplot[color=blue, mark=o, thick, dashed] coordinates {(S1,670) (S2,1235) (S3,1985) (S4,2835)};
\addlegendentry{PV Cumulativo}
\addplot[color=teal, mark=x, thick] coordinates {(S1,670) (S2,1235) (S3,1985) (S4,2835)};
\addlegendentry{EV Cumulativo}
\end{axis}
\end{tikzpicture}
\caption{Grafico dell'andamento cumulativo di Planned Value ed Earned Value per sprint.}
\end{figure}

\noindent\textbf{Valutazione:} Entrambi i valori sono $\geq 0$ per tutti gli \glslink{sprint}{sprint} (\emph{soglia accettabile} raggiunta). Il PV cumulativo è $\leq \text{BAC} = \EUR{12.575,00}$ (\emph{soglia ottimale per MPC-PV} raggiunta); l'EV cumulativo è $\leq \text{EAC}$ (\emph{soglia ottimale per MPC-EV} raggiunta).

% ─────────────────────────────────────────────────────────────────────────────
\subsubsection{MPC-AC — Actual Cost}

L'Actual Cost (AC) rappresenta il costo effettivo sostenuto per il lavoro completato in ciascuno \glslink{sprint}{sprint}, ricavato dai consuntivi del \glslink{piano-di-progetto}{Piano di Progetto}.

\begin{table}[H]
\centering
\rowcolors{2}{gray!15}{white}
\begin{tabular}{|l|r|r|}
    \hline
    \rowcolor{gray!35}
    \textbf{\glslink{sprint}{Sprint}} & \textbf{AC (\euro)} & \textbf{AC cumulativo (\euro)} \\
    \hline
    S1 & 705,00  & 705,00  \\
    S2 & 625,00  & 1.330,00 \\
    S3 & 830,00  & 2.160,00 \\
    S4 & 975,00  & 3.135,00 \\
    \hline
\end{tabular}
\caption{Andamento MPC-AC per sprint (fase RTB)}
\end{table}

\begin{figure}[H]
\centering
\begin{tikzpicture}
\begin{axis}[
    title={Andamento MPC-AC Cumulativo},
    xlabel={\glslink{sprint}{Sprint}},
    ylabel={Costo Effettivo (\euro)},
    symbolic x coords={S1, S2, S3, S4},
    xtick=data,
    ymin=0, ymax=3500,
    grid=both,
    grid style={dashed, gray!30},
    legend pos=north west,
    width=0.85\textwidth,
    height=6.5cm
]
\addplot[color=red!80!black, mark=square*, thick] coordinates {(S1,705) (S2,1330) (S3,2160) (S4,3135)};
\addlegendentry{AC Cumulativo}
\end{axis}
\end{tikzpicture}
\caption{Grafico dell'andamento dei costi effettivi cumulati.}
\end{figure}

\noindent\textbf{Valutazione:} I costi effettivi sono positivi in tutti gli \glslink{sprint}{sprint} (\emph{soglia accettabile} raggiunta). Rispetto al PV si osserva un lieve sforamento in ogni \glslink{sprint}{sprint}, indicando costi leggermente superiori a quelli pianificati.

% ─────────────────────────────────────────────────────────────────────────────
\subsubsection{MPC-CPI e MPC-SPI — Cost Performance Index e Schedule Performance Index}

Il Cost Performance Index misura l'efficienza dei costi:
\[
  \text{CPI} = \frac{\text{EV}_\text{cum}}{\text{AC}_\text{cum}}
\]
Lo Schedule Performance Index misura l'efficienza del tempo:
\[
  \text{SPI} = \frac{\text{EV}_\text{cum}}{\text{PV}_\text{cum}}
\]

\begin{table}[H]
\centering
\rowcolors{2}{gray!15}{white}
\begin{tabular}{|l|r|r|}
    \hline
    \rowcolor{gray!35}
    \textbf{\glslink{sprint}{Sprint}} & \textbf{CPI} & \textbf{SPI} \\
    \hline
    S1 & 0,950 & 1,000 \\
    S2 & 0,929 & 1,000 \\
    S3 & 0,919 & 1,000 \\
    S4 & 0,904 & 1,000 \\
    \hline
\end{tabular}
\caption{Andamento MPC-CPI e MPC-SPI per sprint (fase RTB)}
\end{table}

\begin{figure}[H]
\centering
\begin{tikzpicture}
\begin{axis}[
    title={Andamento indici CPI e SPI},
    xlabel={\glslink{sprint}{Sprint}},
    ylabel={Indice},
    symbolic x coords={S1, S2, S3, S4},
    xtick=data,
    ymin=0.8, ymax=1.1,
    grid=both,
    grid style={dashed, gray!30},
    legend pos=south west,
    width=0.85\textwidth,
    height=6.5cm
]
\addplot[color=orange, mark=violet, thick] coordinates {(S1,0.950) (S2,0.929) (S3,0.919) (S4,0.904)};
\addlegendentry{CPI}
\addplot[color=cyan!80!black, mark=triangle*, thick] coordinates {(S1,1.000) (S2,1.000) (S3,1.000) (S4,1.000)};
\addlegendentry{SPI}
\draw[red, thick, dashed] (axis cs:S1,0.9) -- (axis cs:S4,0.9);
\node[red, anchor=south west] at (axis cs:S1,0.9) {\footnotesize Soglia Accettabile (0.9)};
\end{axis}
\end{tikzpicture}
\caption{Confronto tra Cost Performance Index e Schedule Performance Index.}
\end{figure}

\noindent\textbf{Valutazione:} Il CPI risulta in lieve calo progressivo ma rimane $\geq 0{,}9$ (\emph{soglia accettabile} raggiunta) per tutti gli \glslink{sprint}{sprint}. Il valore a S4 di 0,904 indica che per ogni euro pianificato si è sostenuto un costo di circa \EUR{1,11}. L'SPI pari a 1,000 in tutti gli \glslink{sprint}{sprint} conferma il pieno rispetto delle scadenze pianificate (\emph{soglia ottimale} raggiunta).

% ─────────────────────────────────────────────────────────────────────────────
\subsubsection{MPC-EAC e MPC-ETC — Estimated At Completion e Estimate To Complete}

L'Estimated At Completion stima il costo totale del progetto proiettando la performance corrente:
\[
  \text{EAC} = \frac{\text{BAC}}{\text{CPI}}
\]
L'Estimate To Complete rappresenta la stima dei costi rimanenti al completamento:
\[
  \text{ETC} = \text{EAC} - \text{AC}_\text{cum}
\]

\begin{table}[H]
\centering
\rowcolors{2}{gray!15}{white}
\begin{tabular}{|l|r|r|}
    \hline
    \rowcolor{gray!35}
    \textbf{\glslink{sprint}{Sprint}} & \textbf{EAC (\euro)} & \textbf{ETC (\euro)} \\
    \hline
    S1 & 13.231,00 & 12.526,00 \\
    S2 & 13.544,00 & 12.214,00 \\
    S3 & 13.683,00 & 11.523,00 \\
    S4 & 13.906,00 & 10.771,00 \\
    \hline
\end{tabular}
\caption{Andamento MPC-EAC e MPC-ETC per sprint (fase RTB)}
\end{table}

\begin{figure}[H]
\centering
\begin{tikzpicture}
\begin{axis}[
    title={Andamento EAC e ETC rispetto al BAC},
    xlabel={\glslink{sprint}{Sprint}},
    ylabel={Valore (\euro)},
    symbolic x coords={S1, S2, S3, S4},
    xtick=data,
    ymin=9000, ymax=15000,
    grid=both,
    grid style={dashed, gray!30},
    legend pos=north east,
    width=0.85\textwidth,
    height=6.5cm
]
\addplot[color=blue, mark=purple, thick] coordinates {(S1,13231) (S2,13544) (S3,13683) (S4,13906)};
\addlegendentry{EAC}
\addplot[color=olive, mark=diamond*, thick] coordinates {(S1,12526) (S2,12214) (S3,11523) (S4,10771)};
\addlegendentry{ETC}
\draw[green!60!black, ultra thick, dotted] (axis cs:S1,12575) -- (axis cs:S4,12575);
\node[green!60!black, anchor=north west] at (axis cs:S1,12575) {\footnotesize BAC (12.575 \euro)};
\end{axis}
\end{tikzpicture}
\caption{Proiezione dei costi finali (EAC) e costi a finire (ETC) rispetto al budget.}
\end{figure}

\noindent\textbf{Valutazione:} L'EAC stimato supera il BAC di \EUR{12.575,00} in tutti gli \glslink{sprint}{sprint} e si attesta a \EUR{13.906,00} a S4, con uno scostamento del $+10{,}6\%$ rispetto al BAC — superiore alla soglia accettabile di $\pm 5\%$ (\emph{soglia non raggiunta} per MPC-EAC). L'ETC è positivo per tutti gli \glslink{sprint}{sprint} (\emph{soglia accettabile} raggiunta per MPC-ETC). Il trend crescente dell'EAC richiede un monitoraggio rigoroso dei costi nelle fasi successive.

% ─────────────────────────────────────────────────────────────────────────────
\subsubsection{MPC-TEAC — Time Estimate At Completion}

Il Time Estimate At Completion stima la durata totale del progetto a partire dalla performance di avanzamento. Calcolato come $\text{TEAC} = \text{Durata pianificata} / \text{SPI}$: poiché $\text{SPI} = 1{,}000$ per tutti gli \glslink{sprint}{sprint}, il progetto procede esattamente in linea con il calendario.

\begin{table}[H]
\centering
\rowcolors{2}{gray!15}{white}
\begin{tabular}{|l|r|r|}
    \hline
    \rowcolor{gray!35}
    \textbf{\glslink{sprint}{Sprint}} & \textbf{SPI} & \textbf{TEAC (settimane)} \\
    \hline
    S1 & 1,000 & 8,0 \\
    S2 & 1,000 & 8,0 \\
    S3 & 1,000 & 8,0 \\
    S4 & 1,000 & 8,0 \\
    \hline
\end{tabular}
\caption{Andamento MPC-TEAC per sprint (fase RTB — durata pianificata RTB: 8 settimane)}
\end{table}

\begin{figure}[H]
\centering
\begin{tikzpicture}
\begin{axis}[
    title={Andamento Time Estimate At Completion},
    xlabel={\glslink{sprint}{Sprint}},
    ylabel={Settimane},
    symbolic x coords={S1, S2, S3, S4},
    xtick=data,
    ymin=0, ymax=10,
    grid=both,
    grid style={dashed, gray!30},
    legend pos=north west,
    width=0.85\textwidth,
    height=6.5cm
]
\addplot[color=violet, mark=*, thick] coordinates {(S1,8.0) (S2,8.0) (S3,8.0) (S4,8.0)};
\addlegendentry{TEAC}
\draw[gray, thick, dashed] (axis cs:S1,8.0) -- (axis cs:S4,8.0);
\node[gray, anchor=south west] at (axis cs:S1,8.0) {\footnotesize Durata Pianificata (8 sett.)};
\end{axis}
\end{tikzpicture}
\caption{Stima del tempo finale di conclusione delle attività.}
\end{figure}

\noindent\textbf{Valutazione:} TEAC = durata pianificata (8 settimane) in tutti gli \glslink{sprint}{sprint} (\emph{soglia ottimale} raggiunta). Il progetto non presenta ritardi di \glslink{pianificazione}{pianificazione} nella fase \glslink{rtb-requirements-and-technology-baseline}{RTB}.

% ─────────────────────────────────────────────────────────────────────────────
\subsection{Processi primari — Sviluppo}

\subsubsection{MPC-RSI — Requirements Stability Index}

Il Requirements Stability Index valuta la stabilità dei requisiti nel tempo. La baseline dei requisiti è stata consolidata al termine dello \glslink{sprint}{Sprint} 3 con un totale di 131 requisiti (109 funzionali, 6 di qualità, 16 di vincolo). Durante lo \glslink{sprint}{Sprint} 4 sono state apportate correzioni ad alcuni requisiti esistenti.

\begin{table}[H]
    \centering
    \rowcolors{2}{gray!15}{white}
    \begin{tabular}{|l|c|c|c|c|c|c|}
        \hline
        \rowcolor{gray!35}
        \textbf{\glslink{sprint}{Sprint}} & \textbf{Baseline} & \textbf{Aggiunti} & \textbf{Modificati} & \textbf{Eliminati} & \textbf{RSI (\%)} & \textbf{Esito} \\
        \hline
        S1 & N/A & N/A & N/A & N/A & N/A & N/A \\
        S2 & N/A & N/A & N/A & N/A & N/A & N/A \\
        S3 & 131 & 0 & 0 & 0 & 100\% & \textcolor{green!60!black}{Ottimo} \\
        S4 & 131 & 0 & 4 & 0 & 96{,}9\% & \textcolor{green!60!black}{Accettabile} \\
        \hline
    \end{tabular}
    \caption{Andamento MPC-RSI per sprint}
\end{table}

\begin{figure}[H]
\centering
\begin{tikzpicture}
\begin{axis}[
    title={Andamento Requirements Stability Index},
    xlabel={\glslink{sprint}{Sprint}},
    ylabel={RSI (\%)},
    symbolic x coords={S1, S2, S3, S4},
    xtick=data,
    ymin=70, ymax=105,
    grid=both,
    grid style={dashed, gray!30},
    width=0.85\textwidth,
    height=6.5cm
]
\addplot[color=teal, mark=*, thick] coordinates {(S3,100) (S4,96.9)};
\addlegendentry{RSI Rilevato}
\draw[red, thick, dashed] (axis cs:S1,70) -- (axis cs:S4,70);
\node[red, anchor=south west] at (axis cs:S1,70) {\footnotesize Soglia Accettabile (70\%)};
\end{axis}
\end{tikzpicture}
\caption{Grafico della stabilità dei requisiti post-consolidamento.}
\end{figure}

\noindent\textbf{Valutazione:} La misurazione inizia formalmente dallo \glslink{sprint}{Sprint} 3, momento di consolidamento della baseline. Nello \glslink{sprint}{Sprint} 4 sono state effettuate 4 correzioni formali. Il valore si attesta al 96,9\%, mantenendosi ampiamente sopra la soglia accettabile del 70\% (\emph{soglia accettabile} raggiunta).

% ─────────────────────────────────────────────────────────────────────────────
\subsection{Processi di supporto — Documentazione}

\subsubsection{MPC-IG — Indice di Gulpease}

L'\glslink{indice-di-gulpease-mpc-ig}{Indice di Gulpease} misura la leggibilità dei documenti in lingua italiana secondo la formula:
\[
  G = 89 - \frac{10 \cdot L}{P} + \frac{300 \cdot F}{P}
\]
dove $L$ è il numero di lettere, $P$ il numero di parole e $F$ il numero di frasi. Il calcolo è stato effettuato tramite estrazione del testo con \texttt{\glslink{detex}{detex}} e script Python.

\begin{table}[H]
\centering
\rowcolors{2}{gray!15}{white}
\begin{tabular}{|l|r|c|c|}
    \hline
    \rowcolor{gray!35}
    \textbf{Documento} & \textbf{\glslink{indice-di-gulpease-mpc-ig}{Indice di Gulpease}} & \textbf{Soglia accettabile} & \textbf{Esito} \\
    \hline
    \glslink{analisi-dei-requisiti-adr}{Analisi dei Requisiti} & 59,7 & $\geq 60$ & Borderline \\
    \glslink{norme-di-progetto}{Norme di Progetto}     & 45,9 & $\geq 60$ & \textcolor{red!70!black}{Non soddisfatta} \\
    \glslink{piano-di-progetto}{Piano di Progetto}     & 47,7 & $\geq 60$ & \textcolor{red!70!black}{Non soddisfatta} \\
    \glslink{piano-di-qualifica}{Piano di Qualifica}    & 55,8 & $\geq 60$ & \textcolor{red!70!black}{Non soddisfatta} \\
    \hline
\end{tabular}
\caption{Indice di Gulpease dei documenti RTB}
\end{table}

\begin{figure}[H]
\centering
\begin{tikzpicture}
\begin{axis}[
    title={\glslink{indice-di-gulpease-mpc-ig}{Indice di Gulpease} per Documento},
    ylabel={Punteggio Indice},
    symbolic x coords={\glslink{analisi-dei-requisiti-adr}{AdR}, NdP, PdP, PdQ},
    xtick=data,
    ymin=0, ymax=80,
    ybar,
    bar width=20pt,
    grid=major,
    grid style={dashed, gray!30},
    width=0.85\textwidth,
    height=6.5cm,
    nodes near coords,
    nodes near coords align={vertical}
]
\addplot[fill=blue!35!white, draw=black] coordinates {
    (\glslink{analisi-dei-requisiti-adr}{AdR},59.7) 
    (NdP,45.9) 
    (PdP,47.7) 
    (PdQ,55.8)
};
\draw[red, thick, dashed] (axis cs:AdR,60) -- (axis cs:PdQ,60);
\node[red, anchor=south east] at (axis cs:PdQ,60) {\footnotesize Soglia Accettabile (60)};
\end{axis}
\end{tikzpicture}
\caption{Livello di leggibilità dei documenti calcolato tramite Indice di Gulpease.}
\end{figure}

\noindent\textbf{Valutazione:} I documenti tecnici mostrano punteggi compresi tra 45,9 e 59,7, al di sotto della soglia accettabile di 60 (\emph{soglia non raggiunta}). Questo risultato è strutturalmente atteso: la terminologia tecnica specialistica, gli acronimi e i costrutti formali riducono l'indice di leggibilità nei documenti ingegneristici. Si raccomanda di semplificare ove possibile la struttura delle frasi nelle prossime revisioni.

% ─────────────────────────────────────────────────────────────────────────────
\subsubsection{MPC-CO — Correttezza ortografica}

La verifica è stata condotta tramite estrazione del testo con \texttt{\glslink{detex}{detex}} e analisi automatica della correttezza ortografica.

\begin{table}[H]
\centering
\rowcolors{2}{gray!15}{white}
\begin{tabular}{|l|r|c|}
    \hline
    \rowcolor{gray!35}
    \textbf{Documento} & \textbf{Errori ortografici} & \textbf{Esito} \\
    \hline
    \glslink{analisi-dei-requisiti-adr}{Analisi dei Requisiti} & 0 & \textcolor{green!60!black}{Soddisfatta} \\
    \glslink{norme-di-progetto}{Norme di Progetto}     & 0 & \textcolor{green!60!black}{Soddisfatta} \\
    \glslink{piano-di-progetto}{Piano di Progetto}     & 0 & \textcolor{green!60!black}{Soddisfatta} \\
    \glslink{piano-di-qualifica}{Piano di Qualifica}    & 0 & \textcolor{green!60!black}{Soddisfatta} \\
    \hline
\end{tabular}
\caption{Correttezza ortografica dei documenti RTB}
\end{table}

\noindent\textbf{Valutazione:} Tutti i documenti presentano 0 errori ortografici (\emph{soglia ottimale} raggiunta).

% ─────────────────────────────────────────────────────────────────────────────
\subsection{Processi di supporto — Verifica}

\subsubsection{MPC-CC e MPC-TSR — Code Coverage e Test Success Rate}

Le metriche MPC-CC (Code Coverage) e MPC-TSR (Test Success Rate) non sono ancora considerate pienamente applicabili nella fase \glslink{rtb-requirements-and-technology-baseline}{RTB}, poiché il codice realizzato in \glslink{rtb-requirements-and-technology-baseline}{RTB} ha natura dimostrativa e non coincide con una baseline stabile di prodotto. La misurazione sistematica di tali metriche è prevista a partire dalla fase \glslink{pb-product-baseline}{PB (Product Baseline)}, quando verrà sviluppata e consolidata la baseline di riferimento e i test automatici potranno essere utilizzati come criterio continuativo di verifica.

% ─────────────────────────────────────────────────────────────────────────────
\subsection{Processi di supporto — Gestione della qualità}

\subsubsection{MPC-QMS — Quality Metrics Satisfied}

Il Quality Metrics Satisfied misura la percentuale di metriche di processo che soddisfano la soglia accettabile nel periodo di riferimento:
\[
  \text{QMS} = \frac{\text{Metriche soddisfatte}}{\text{Metriche misurate}} \times 100\%
\]

La tabella seguente riassume l'esito di tutte le metriche misurabili nella fase \glslink{rtb-requirements-and-technology-baseline}{RTB}. Con l'inclusione delle metriche MPC-RSI e MPC-PRI, le metriche misurate passano a 12.

\begin{table}[H]
\centering
\rowcolors{2}{gray!15}{white}
\begin{tabular}{|l|p{3.5cm}|p{3cm}|c|}
    \hline
    \rowcolor{gray!35}
    \textbf{Metrica} & \textbf{Valore misurato} & \textbf{Soglia accettabile} & \textbf{Esito} \\
    \hline
    MPC-EV  & $\geq 0$ tutti gli \glslink{sprint}{sprint} & $\geq 0$ & \textcolor{green!60!black}{Soddisfatta} \\
    MPC-PV  & $\geq 0$ tutti gli \glslink{sprint}{sprint} & $\geq 0$ & \textcolor{green!60!black}{Soddisfatta} \\
    MPC-AC  & $\geq 0$ tutti gli \glslink{sprint}{sprint} & $\geq 0$ & \textcolor{green!60!black}{Soddisfatta} \\
    MPC-CPI & 0,904 (S4)               & $\geq 0{,}9$ & \textcolor{green!60!black}{Soddisfatta} \\
    MPC-SPI & 1,000 (tutti gli \glslink{sprint}{sprint}) & $\geq 0{,}9$ & \textcolor{green!60!black}{Soddisfatta} \\
    MPC-EAC & $+10{,}6\%$ dal BAC (S4) & $\pm 5\%$ dal BAC & \textcolor{red!70!black}{Non soddisfatta} \\
    MPC-ETC & $\geq 0$ tutti gli \glslink{sprint}{sprint} & $\geq 0$ & \textcolor{green!60!black}{Soddisfatta} \\
    MPC-TEAC & 8 settimane (SPI = 1,000) & $\geq 0$ & \textcolor{green!60!black}{Soddisfatta} \\
    \glslink{indice-di-gulpease-mpc-ig}{MPC-IG}  & 45,9 – 59,7             & $\geq 60$ & \textcolor{red!70!black}{Non soddisfatta} \\
    MPC-CO  & 0 errori                 & 0 errori & \textcolor{green!60!black}{Soddisfatta} \\
    MPC-RSI & 96,9\% (S4)              & $\geq 70\%$ & \textcolor{green!60!black}{Soddisfatta} \\
    MPC-PRI & 50\% (fase \glslink{rtb-requirements-and-technology-baseline}{RTB})          & $\leq 20\%$ & \textcolor{red!70!black}{Non soddisfatta} \\
    \hline
    \rowcolor{gray!25}
    \multicolumn{3}{|l|}{\textbf{Metriche soddisfatte / Metriche misurate}} & \textbf{9 / 12} \\
    \hline
    \rowcolor{gray!25}
    \multicolumn{3}{|l|}{\textbf{MPC-QMS}} & \textbf{75\%} \\
    \hline
\end{tabular}
\caption{Calcolo MPC-QMS per la fase RTB aggiornato con RSI e PRI}
\end{table}

L'indice MPC-QMS riflette la conformità complessiva dei processi del team. In S1, la soddisfazione è massima (100\%) per i dati valutati in quel frangente. Nei successivi S2 e S3 il valore si stabilizza all'80\%. Nello \glslink{sprint}{Sprint} 4, con l'inclusione delle metriche globali RSI e PRI del periodo \glslink{rtb-requirements-and-technology-baseline}{RTB}, il QMS scende al 75\% (essendo MPC-PRI non soddisfatta). Il grafico seguente mostra questo andamento, confrontato con le soglie di accettabilità e ottimalità.

\begin{figure}[H]
\centering
\begin{tikzpicture}
\begin{axis} [
    title={Andamento MPC-QMS per \glslink{sprint}{Sprint} rispetto alle Soglie},
    ylabel={Percentuale soddisfazione (\%)},
    symbolic x coords={S1, S2, S3, S4},
    xtick=data,
    ymin=0, ymax=100,
    ytick={0, 20, 40, 60, 80, 100},
    grid=major,
    grid style={dashed, gray!30},
    bar width=18pt, 
    legend style={at={(0.5,-0.20)}, anchor=north, legend columns=-1, draw=none} 
]

\addplot[ybar stacked, fill=gray!8, pattern=crosshatch, pattern color=black!40, forget plot, bar width=24pt] coordinates {(S1,20) (S2,20) (S3,20) (S4,20)};

\addplot[ybar stacked, fill=none, pattern=dots, pattern color=black!40, forget plot, bar width=24pt] coordinates {(S1,80) (S2,80) (S3,80) (S4,80)};

\addplot[ybar, fill=blue!70!white, draw=none, bar width=10pt, forget plot] coordinates {(S1,100) (S2,0) (S3,0) (S4,0)};

\addplot[ybar, fill=green!60!black, draw=none, bar width=10pt, forget plot] coordinates {(S1,0) (S2,80) (S3,80) (S4,0)};

\addplot[ybar, fill=orange!80!white, draw=none, bar width=10pt, forget plot] coordinates {(S1,0) (S2,0) (S3,0) (S4,75)};

\addlegendimage{ybar, fill=blue!70!white, draw=none}
\addlegendentry{Valore Rilevato}

\addlegendimage{area legend, pattern=dots, pattern color=black!40, fill=none}
\addlegendentry{Valore Accettabile (Soglia 80\%)}

\addlegendimage{area legend, fill=gray!8, pattern=crosshatch, pattern color=black!40}
\addlegendentry{Valore Ottimale (Soglia 100\%)}

\end{axis}
\end{tikzpicture}
\caption{Grafico percentuale di metriche soddisfatte (MPC-QMS) per sprint confrontato con le soglie di riferimento.}
\end{figure}

\noindent\textbf{Valutazione:} Il QMS complessivo al termine della fase \glslink{rtb-requirements-and-technology-baseline}{RTB} è pari al 75\%, al di sotto della \emph{soglia accettabile} ($\geq 80\%$). Le tre metriche non soddisfatte richiedono interventi correttivi:
\begin{itemize}
    \item \textbf{MPC-EAC:} il costo stimato a completamento supera il BAC del 10,6\%. Nelle fasi successive è necessario un controllo più rigoroso dei costi, con revisione periodica delle stime e delle allocazioni orarie.
    \item \textbf{MPC-PRI:} il verificarsi di imprevisti bloccanti (50\% rispetto \glslink{intelligenza-artificiale-ai}{ai} rischi emersi) suggerisce che il registro \glslink{mitigazione-dei-rischi}{dei rischi} necessiti di un ampliamento e aggiornamento in vista dello sviluppo operativo \glslink{pb-product-baseline}{PB}.
    \item \textbf{\glslink{indice-di-gulpease-mpc-ig}{MPC-IG}:} i documenti tecnici presentano indici di leggibilità al di sotto della soglia di 60. Si raccomanda una revisione stilistica nelle prossime versioni, semplificando la struttura delle frasi.
\end{itemize}

% ─────────────────────────────────────────────────────────────────────────────
\subsection{Processi organizzativi — Gestione dei processi}

\subsubsection{MPC-TE — Time Efficiency}

Durante la fase \glslink{rtb-requirements-and-technology-baseline}{RTB}, il tracciamento orario si è concentrato esclusivamente sullo scostamento tra le ore previste e le ore effettivamente impiegate per il completamento dei task. Non è stata effettuata una misurazione granulare per distinguere le ore produttive da quelle di \textit{overhead}. Pertanto, questa metrica risulta \textbf{N/A} per gli \glslink{sprint}{sprint} S1-S4.

\noindent\textbf{Azioni di miglioramento:} Il team introdurrà un sistema di \textit{time-tracking} granulare a partire dalla fase di \glslink{pb-product-baseline}{Product Baseline} (\glslink{pb-product-baseline}{PB}).

% ─────────────────────────────────────────────────────────────────────────────
\subsubsection{MPC-PRI — Percentuale Rischi Inattesi}

La metrica valuta la capacità del team di prevedere e mitigare i rischi. I dati sono stati estratti dalle retrospettive dei quattro \glslink{sprint}{sprint} di progetto.

\begin{table}[H]
    \centering
    \renewcommand{\arraystretch}{1.2}
    \rowcolors{2}{gray!15}{white}
    \begin{tabularx}{\textwidth}{|l|X|c|}
        \hline
        \rowcolor{gray!35}
        \textbf{\glslink{sprint}{Sprint}} & \textbf{Problema Emerso} & \textbf{Classificazione} \\
        \hline
        S1 & Difficoltà nell'apprendimento di nuovi tool grafici per il mock-up. & Previsto (RT1) \\
        \hline
        S2 & Stallo per attesa file di esempio dal \glslink{committente-proponente}{proponente}. & Previsto (RR1 / RO5) \\
        \hline
        S3 & Sottostima del lavoro per diagrammazione e attività tecniche non preventivate. & Inatteso \\
        \hline
        S4 & Blocco per attesa credenziali infrastruttura cloud AWS. & Inatteso \\
        \hline
    \end{tabularx}
    \caption{Tracciamento dei rischi manifestati per sprint}
\end{table}

\begin{table}[H]
    \centering
    \rowcolors{2}{gray!15}{white}
    \begin{tabular}{|l|c|c|c|c|}
        \hline
        \rowcolor{gray!35}
        \textbf{\glslink{sprint}{Sprint}} & \textbf{Rischi Totali} & \textbf{Rischi Inattesi} & \textbf{PRI (\%)} & \textbf{Esito} \\
        \hline
        Fase \glslink{rtb-requirements-and-technology-baseline}{RTB} & 4 & 2 & 50\% & \textcolor{red!70!black}{Non soddisfatta} \\
        \hline
    \end{tabular}
    \caption{Calcolo MPC-PRI per la fase RTB}
\end{table}

\noindent\textbf{Valutazione:} Durante la fase \glslink{rtb-requirements-and-technology-baseline}{RTB} sono emersi 4 problemi principali. Due di questi non erano stati esplicitamente previsti. Il valore PRI si attesta al 50\%, superando la soglia accettabile del $\leq 20\%$ (\emph{soglia non raggiunta}). Sarà necessaria un'integrazione del registro \glslink{mitigazione-dei-rischi}{dei rischi} per la fase \glslink{pb-product-baseline}{PB}.

% ─────────────────────────────────────────────────────────────────────────────
\subsubsection{MPC-LE — Labor Efficiency}

Come per la metrica MPC-TE, l'assenza di un tracciamento puntuale delle ore spese in attività a valore aggiunto rispetto al tempo totale lavorato (incluso overhead) rende impossibile il calcolo per la fase \glslink{rtb-requirements-and-technology-baseline}{RTB}. Il valore è quindi \textbf{N/A}. Le misurazioni inizieranno formalmente con la fase \glslink{pb-product-baseline}{PB}.

\section{Iniziative di miglioramento}

Le valutazioni riportate nel cruscotto di qualità evidenziano alcuni ambiti sui quali il gruppo intende intervenire nella fase \glslink{pb-product-baseline}{PB}. Il valore finale di MPC-QMS è pari al 75\%, inferiore alla soglia accettabile, e segnala la necessità di trasformare le criticità rilevate in azioni operative verificabili. Gli interventi previsti non introducono nuove metriche, ma utilizzano quelle già definite nel presente documento come riferimento per controllare l'efficacia delle correzioni.

\begin{table}[H]
    \centering
    \renewcommand{\arraystretch}{1.25}
    \rowcolors{2}{gray!15}{white}
    \begin{tabularx}{\textwidth}{|p{3.1cm}|X|X|}
        \hline
        \rowcolor{gray!35}
        \textbf{Ambito} & \textbf{Evidenza rilevata in \glslink{rtb-requirements-and-technology-baseline}{RTB}} & \textbf{Obiettivo per la fase \glslink{pb-product-baseline}{PB}} \\
        \hline
        \glslink{pianificazione}{Pianificazione} & MPC-EAC non soddisfatta, con EAC superiore al BAC del 10,6\%. & Raffinare le stime e controllare con maggiore frequenza le previsioni a finire. \\
        \hline
        Gestione \glslink{mitigazione-dei-rischi}{dei rischi} & MPC-PRI non soddisfatta, con 2 rischi inattesi su 4 e valore pari al 50\%. & Rendere più sistematica l'identificazione \glslink{mitigazione-dei-rischi}{dei rischi} tecnici, organizzativi ed esterni. \\
        \hline
        Qualità documentale & \glslink{indice-di-gulpease-mpc-ig}{MPC-IG} sotto soglia per la leggibilità dei documenti. & Ridurre la densità testuale e migliorare la distinzione tra definizioni, motivazioni e conseguenze. \\
        \hline
        Tracciamento operativo & MPC-TE e MPC-LE non misurabili in \glslink{rtb-requirements-and-technology-baseline}{RTB} per assenza di tracciamento granulare. & Distinguere le ore per tipologia di attività, ruolo, \glslink{sprint}{sprint} e \glslink{issue}{issue} di riferimento. \\
        \hline
        Verifica e testing & MPC-CC e MPC-TSR non ancora misurabili su una baseline stabile di prodotto, da sviluppare in \glslink{pb-product-baseline}{PB}. & Introdurre quality gate e tracciamento continuativo dei test a partire dai componenti stabilizzati. \\
        \hline
    \end{tabularx}
    \caption{Sintesi delle criticità RTB e degli obiettivi di miglioramento per PB}
\end{table}

Le azioni descritte nelle sottosezioni seguenti sono state definite a partire dagli esiti delle metriche e dalle evidenze emerse durante gli \glslink{sprint}{sprint} \glslink{rtb-requirements-and-technology-baseline}{RTB}. Per ciascun ambito sono indicati la criticità rilevata, l'intervento previsto e il criterio con cui il gruppo intende verificarne l'efficacia nella fase successiva.

\subsection{Miglioramento della pianificazione e del controllo dei costi}

Lo scostamento rilevato tramite MPC-EAC evidenzia che alcune attività sono state stimate con un livello di dettaglio non sufficiente rispetto all'impegno effettivo richiesto. La criticità riguarda in particolare le attività tecniche e documentali, che nella fase \glslink{rtb-requirements-and-technology-baseline}{RTB} hanno richiesto più tempo rispetto a quanto preventivato.

\begin{table}[H]
    \centering
    \renewcommand{\arraystretch}{1.25}
    \rowcolors{2}{gray!15}{white}
    \begin{tabularx}{\textwidth}{|p{3cm}|X|X|}
        \hline
        \rowcolor{gray!35}
        \textbf{Aspetto} & \textbf{Criticità rilevata} & \textbf{Intervento previsto in \glslink{pb-product-baseline}{PB}} \\
        \hline
        Stima delle attività & Alcune attività sono state valutate in modo troppo aggregato, rendendo meno visibili dipendenze e sottoattività. & Scomporre le attività più ampie in task verificabili e associare a ciascuno una stima oraria più esplicita. \\
        \hline
        Controllo dello \glslink{sprint}{sprint} & Lo scostamento è emerso soprattutto a \glslink{consuntivo}{consuntivo}, quando l'impatto sul costo stimato a completamento era già visibile. & Confrontare a fine \glslink{sprint}{sprint} ore preventivate e consuntivate, motivando gli scostamenti principali. \\
        \hline
        Previsione a finire & L'EAC ha superato il BAC del 10,6\%, oltre la soglia accettabile prevista. & Aggiornare con maggiore frequenza le previsioni a finire, tenendo conto dei dati consuntivi più recenti. \\
        \hline
    \end{tabularx}
    \caption{Azioni di miglioramento per MPC-EAC}
\end{table}

Per evitare che il controllo dei costi resti limitato al solo \glslink{consuntivo}{consuntivo}, il gruppo intende applicare un controllo regolare sulle stime ancora aperte. Tale controllo sarà svolto in corrispondenza della chiusura degli \glslink{sprint}{sprint} e dovrà rendere esplicite le cause dello scostamento, distinguendo tra sottostima iniziale, attività non previste, dipendenze esterne e riallocazione delle risorse.

\paragraph*{Criticità rilevata:} MPC-EAC non soddisfatta, con EAC superiore al BAC del 10,6\%.

\paragraph*{Azione prevista:} Confronto sistematico tra ore preventivate e ore consuntivate, con aggiornamento delle previsioni a finire sulla base degli scostamenti rilevati.

\paragraph*{Metriche monitorate:} MPC-EAC, con supporto delle metriche MPC-AC, MPC-CPI e MPC-ETC.

\subsection{Miglioramento della gestione dei rischi}

Il valore di MPC-PRI mostra una capacità di previsione \glslink{mitigazione-dei-rischi}{dei rischi} ancora incompleta. In \glslink{rtb-requirements-and-technology-baseline}{RTB} alcuni eventi sono stati gestiti nel corso dello \glslink{sprint}{sprint}, ma non erano stati descritti in modo sufficientemente esplicito nel registro \glslink{mitigazione-dei-rischi}{dei rischi}. Per la fase \glslink{pb-product-baseline}{PB} il gruppo intende quindi rendere più puntuale la classificazione \glslink{mitigazione-dei-rischi}{dei rischi} e collegarla alle retrospettive.

\begin{table}[H]
    \centering
    \renewcommand{\arraystretch}{1.25}
    \rowcolors{2}{gray!15}{white}
    \begin{tabularx}{\textwidth}{|p{3.2cm}|X|X|}
        \hline
        \rowcolor{gray!35}
        \textbf{\glslink{rischio}{Rischio} emerso} & \textbf{Lettura del problema} & \textbf{Trattamento previsto} \\
        \hline
        Attività tecniche sottostimate & Alcune attività implementative e di integrazione hanno richiesto approfondimenti non completamente previsti. & Esplicitare nel registro i rischi legati alla complessità tecnica e indicare una \glslink{mitigazione-dei-rischi}{mitigazione} tramite spike o task preliminari. \\
        \hline
        Attese per credenziali & L'accesso a risorse esterne può incidere direttamente sulla continuità delle attività tecniche. & Registrare le dipendenze da credenziali e ambienti esterni, indicando \glslink{responsabile}{responsabile} del sollecito e data limite interna. \\
        \hline
        Dipendenze esterne & File, servizi o informazioni provenienti dal \glslink{committente-proponente}{proponente} possono introdurre ritardi non controllabili dal gruppo. & Associare alle dipendenze esterne un impatto previsto e una possibile attività alternativa in caso di blocco. \\
        \hline
    \end{tabularx}
    \caption{Rischi da esplicitare nel registro per la fase PB}
\end{table}

L'aggiornamento del registro non sarà limitato all'aggiunta di nuove voci, ma dovrà migliorare la qualità informativa \glslink{mitigazione-dei-rischi}{dei rischi} già presenti. In particolare, durante le retrospettive il gruppo dovrà verificare se gli eventi emersi nello \glslink{sprint}{sprint} erano stati previsti, se la \glslink{mitigazione-dei-rischi}{mitigazione} era applicabile e se l'impatto rilevato corrispondeva a quello stimato.

\paragraph*{Criticità rilevata:} MPC-PRI non soddisfatta, con presenza di rischi inattesi durante la fase \glslink{rtb-requirements-and-technology-baseline}{RTB}.

\paragraph*{Azione prevista:} Aggiornamento del registro rischi durante le retrospettive, indicando causa, impatto, \glslink{mitigazione-dei-rischi}{mitigazione} e \glslink{responsabile}{responsabile} per ciascun \glslink{rischio}{rischio} significativo.

\paragraph*{Metriche monitorate:} MPC-PRI, con attenzione alla riduzione \glslink{mitigazione-dei-rischi}{dei rischi} inattesi e alla tempestività delle azioni di \glslink{mitigazione-dei-rischi}{mitigazione}.

\subsection{Miglioramento della qualità documentale}

L'\glslink{indice-di-gulpease-mpc-ig}{indice di Gulpease} sotto soglia non implica necessariamente la presenza di documenti errati, ma segnala una leggibilità inferiore a quella attesa. La criticità principale riguarda la densità dei periodi e l'uso frequente di terminologia tecnica, specialmente nelle sezioni in cui vengono motivate scelte architetturali o descritte metriche di controllo.

\begin{table}[H]
    \centering
    \renewcommand{\arraystretch}{1.25}
    \rowcolors{2}{gray!15}{white}
    \begin{tabularx}{\textwidth}{|p{3.2cm}|X|X|}
        \hline
        \rowcolor{gray!35}
        \textbf{Aspetto documentale} & \textbf{Criticità rilevata} & \textbf{Criterio di revisione} \\
        \hline
        Struttura dei periodi & Frasi lunghe e ricche di incisi riducono la chiarezza, pur mantenendo corretto il contenuto. & Semplificare i periodi più densi e separare le informazioni che hanno funzione diversa. \\
        \hline
        Terminologia tecnica & Alcuni concetti vengono spiegati più volte all'interno dei documenti, aumentando la complessità di lettura. & Usare il \glslink{glossario}{Glossario} per alleggerire le spiegazioni ripetute e mantenere nei documenti solo il contesto necessario. \\
        \hline
        Organizzazione logica & Definizioni, motivazioni e conseguenze talvolta compaiono nello stesso paragrafo. & Distinguere meglio le parti descrittive da quelle valutative, così da rendere più immediata la consultazione. \\
        \hline
    \end{tabularx}
    \caption{Azioni di miglioramento per la qualità documentale}
\end{table}

Il miglioramento della leggibilità dovrà mantenere invariato il livello formale dei documenti. L'obiettivo non è ridurre il contenuto tecnico, ma renderlo più consultabile e meno concentrato in paragrafi estesi. Per questo motivo, nelle prossime revisioni sarà utile intervenire soprattutto sulle sezioni più dense, evitando riscritture generalizzate quando il testo risulta già chiaro e coerente.

\paragraph*{Criticità rilevata:} \glslink{indice-di-gulpease-mpc-ig}{MPC-IG} sotto soglia, con \glslink{indice-di-gulpease-mpc-ig}{indice di Gulpease} inferiore al valore accettabile previsto.

\paragraph*{Azione prevista:} Revisione mirata dei periodi più complessi, migliore separazione tra definizioni, motivazioni e conseguenze, uso più sistematico del \glslink{glossario}{Glossario}.

\paragraph*{Metriche monitorate:} \glslink{indice-di-gulpease-mpc-ig}{MPC-IG}, verificando l'andamento dell'indice nelle revisioni successive dei documenti principali.

\subsection{Miglioramento del tracciamento operativo}

Le metriche MPC-TE e MPC-LE non sono state calcolabili in \glslink{rtb-requirements-and-technology-baseline}{RTB} perché il tracciamento orario non distingue ancora in modo sufficiente le diverse categorie di lavoro. Il gruppo ha registrato le ore necessarie al controllo dei costi, ma non ha separato con granularità adeguata attività produttive, verifica, riunioni, studio, overhead e blocchi.

\begin{table}[H]
    \centering
    \renewcommand{\arraystretch}{1.25}
    \rowcolors{2}{gray!15}{white}
    \begin{tabularx}{\textwidth}{|p{3.1cm}|X|X|}
        \hline
        \rowcolor{gray!35}
        \textbf{Categoria} & \textbf{Descrizione operativa} & \textbf{Utilità per le metriche} \\
        \hline
        Attività produttive & Sviluppo, documentazione, configurazione e attività direttamente collegate agli obiettivi dello \glslink{sprint}{sprint}. & Permette di stimare le ore a valore aggiunto necessarie per MPC-TE e MPC-LE. \\
        \hline
        Verifica & Revisione dei documenti, controllo del codice, validazione dei test e correzione dei difetti rilevati. & Consente di valutare il peso della verifica rispetto al lavoro complessivo. \\
        \hline
        Riunioni e coordinamento & Incontri interni, retrospettive, \glslink{pianificazione}{pianificazione} e coordinamento operativo tra membri del gruppo. & Rende visibile il tempo di organizzazione necessario al completamento delle attività. \\
        \hline
        Studio e approfondimento & Analisi di tecnologie, librerie, servizi esterni e strumenti necessari alla realizzazione della PoC o del prodotto. & Distingue il tempo di apprendimento dal tempo di produzione diretta. \\
        \hline
        Overhead e blocchi & Attese, problemi di accesso, dipendenze esterne e attività non produttive ma necessarie alla gestione del lavoro. & Permette di individuare le cause che riducono l'efficienza operativa. \\
        \hline
    \end{tabularx}
    \caption{Categorie operative per il tracciamento delle ore in PB}
\end{table}

La rendicontazione prevista per \glslink{pb-product-baseline}{PB} dovrà collegare ogni registrazione allo \glslink{sprint}{sprint}, all'\glslink{issue}{issue} di riferimento, al ruolo coinvolto, alle ore preventivate e alle ore effettive. Questo consentirà di analizzare gli scostamenti non solo a livello complessivo, ma anche per tipologia di attività e per responsabilità operative.

\paragraph*{Criticità rilevata:} MPC-TE e MPC-LE non misurabili in \glslink{rtb-requirements-and-technology-baseline}{RTB} per mancanza di tracciamento granulare.

\paragraph*{Azione prevista:} Introduzione di una rendicontazione collegata a \glslink{sprint}{sprint}, \glslink{issue}{issue}, ruolo, categoria dell'attività, ore preventivate e ore effettive.

\paragraph*{Metriche monitorate:} MPC-TE e MPC-LE, con supporto dei dati consuntivi utilizzati anche per il controllo di MPC-EAC.

\subsection{Miglioramento della verifica e del testing}

Durante la fase \glslink{rtb-requirements-and-technology-baseline}{RTB} molti test risultano ancora NI o non pienamente misurabili, poiché la PoC ha avuto finalità dimostrativa e ha consentito di verificare la fattibilità dei flussi principali. La baseline stabile di prodotto verrà sviluppata e consolidata nella fase \glslink{pb-product-baseline}{PB}, prendendo come riferimento lo stack effettivamente adottato: \glslink{laravel}{Laravel} con viste Blade, CSS e JavaScript custom. Le tecnologie valutate inizialmente ma non adottate nella realizzazione \glslink{rtb-requirements-and-technology-baseline}{RTB} non saranno considerate nel consolidamento previsto per la \glslink{pb-product-baseline}{PB}.

\begin{table}[H]
    \centering
    \renewcommand{\arraystretch}{1.25}
    \rowcolors{2}{gray!15}{white}
    \begin{tabularx}{\textwidth}{|p{3.3cm}|X|X|}
        \hline
        \rowcolor{gray!35}
        \textbf{Flusso principale} & \textbf{Aspetto da consolidare} & \textbf{Controllo previsto} \\
        \hline
        Caricamento PDF & Validazione del formato, gestione degli errori e tracciamento dell'esito del caricamento. & Test automatici sui casi validi e sui casi di errore più significativi. \\
        \hline
        Salvataggio documenti & Persistenza dei dati generati e recupero corretto delle informazioni associate. & Verifica di creazione, aggiornamento e recupero dei documenti salvati. \\
        \hline
        Elaborazione asincrona & Esecuzione tramite worker e coda \glslink{redis}{Redis}, con gestione degli stati intermedi. & Test sui passaggi di stato e sui casi di elaborazione fallita o ritentata. \\
        \hline
        Integrazione con \glslink{minio}{MinIO} & Gestione dei file e coerenza tra \glslink{metadati}{metadati} applicativi e oggetti salvati. & Controlli su caricamento, lettura e assenza di riferimenti non validi. \\
        \hline
        Uso di Bedrock o dei fallback simulati & Continuità del flusso anche in presenza di limiti o indisponibilità del servizio generativo. & Verifica dei percorsi alternativi e della corretta segnalazione degli errori. \\
        \hline
        Generazione dei contenuti & Produzione e salvataggio dei contenuti generati a partire dai dati estratti. & Test sui risultati attesi e sulle condizioni in cui i dati di partenza sono incompleti. \\
        \hline
    \end{tabularx}
    \caption{Flussi da consolidare nel processo di verifica PB}
\end{table}

In \glslink{pb-product-baseline}{PB} il gruppo intende collegare il tracciamento dei test \glslink{intelligenza-artificiale-ai}{ai} quality gate automatici, così che gli esiti non restino informazioni isolate. Il consolidamento dovrà riguardare sia i test che verificano il comportamento applicativo, sia i controlli di integrazione sui servizi previsti dalla baseline di riferimento.

\paragraph*{Criticità rilevata:} MPC-CC e MPC-TSR non ancora pienamente applicabili in \glslink{rtb-requirements-and-technology-baseline}{RTB}, poiché la baseline stabile di prodotto verrà sviluppata e consolidata in \glslink{pb-product-baseline}{PB}.

\paragraph*{Azione prevista:} Consolidamento dei quality gate automatici e tracciamento degli esiti dei test sui flussi principali del prodotto.

\paragraph*{Metriche monitorate:} MPC-CC e MPC-TSR, da misurare in modo sistematico dopo il consolidamento dei componenti principali.

\bigskip

\noindent\textbf{Sintesi conclusiva:} Le iniziative descritte hanno lo scopo di trasformare le criticità emerse durante la fase \glslink{rtb-requirements-and-technology-baseline}{RTB} in azioni operative per la fase \glslink{pb-product-baseline}{PB}. In particolare, il gruppo intende migliorare la precisione della \glslink{pianificazione}{pianificazione}, rendere più sistematica la gestione \glslink{mitigazione-dei-rischi}{dei rischi}, aumentare la qualità della documentazione e consolidare il processo di verifica del codice. Tali interventi saranno monitorati attraverso le metriche già definite nel presente documento, con particolare attenzione a MPC-EAC, MPC-PRI, \glslink{indice-di-gulpease-mpc-ig}{MPC-IG}, MPC-TE, MPC-LE, MPC-CC e MPC-TSR.

\end{document}
